\documentclass[a4paper,12pt,headsepline, hyphens]{report}

%----------------- PDF CONFIG ----------------- %
\pdfinfo{    
     /Title (Analyse sicherheitsrelevanter Informationen in industriellen Netzwerken für den Einsatz von SIEM Systemen) 
     /Subject   (PDF-Thema)    
     /Author  (Niklas Breuer) 
     /Keywords   (SIEM, Industrie 4.0, Data Collection)      
} 

\title{Analyse sicherheitsrelevanter Informationen in industriellen Netzwerken für den Einsatz von SIEM Systemen}
\author{Niklas Breuer}
\date{??.??.201?}



%----------------- PAKETE INKLUDIEREN ----------------- %

\usepackage{geometry} % Packet für Seitenrandabständex und Einstellung für Seitenränder
\usepackage[ngerman]{babel} % deutsche Silbentrennung

\usepackage{booktabs} %entzerrt die Tabellenzeilen und bietet verschieden dicke Unterteilungslinien
\usepackage{longtable} % Tabellen können sich nicht über mehrere Seiten 
\usepackage{graphicx} % kann LaTeX Grafiken einbinden
\usepackage{tabularx}

\usepackage[utf8]{inputenc} % Umlaute unter Mac werden automatisch gesetzt
\usepackage[T1]{fontenc} % Zeichenencoding
\usepackage{lmodern} % typographische Qualität 
\frenchspacing % Schaltet den zusätzlichen Zwischenraum ab
\usepackage{fix-cm}
\usepackage{hyperref} % verwandelt alle Kapitelüberschriften, Verweise aufs Literaturverzeichnis und andere Querverweise in PDF-Hyperlinks
\usepackage{color}
\usepackage{amsmath} %für mathematische Formeln
\usepackage{amssymb} % ermöglicht die Nutzung verschiedener Sonderzeichen und Symbole
\usepackage[]{acronym} %für das Abkürzungsverzeichnis


%\usepackage[nottoc]{tocbibind}

% für Listings
\usepackage{listings}
\lstset{numbers=left, numberstyle=\tiny, numbersep=5pt, stepnumber=4, keywordstyle=\color{black}\bfseries\itshape, stringstyle=\ttfamily,showstringspaces=false,basicstyle=\footnotesize,captionpos=b}
\lstset{language=java}

%Quellen
\usepackage[numbers]{natbib}
\usepackage{url}

%Farben
\usepackage{color}

%ToDos
\usepackage[colorinlistoftodos]{todonotes}
\newcommand{\todoForm}[1]{\todo[color=yellow!50!black!20,bordercolor=yellow!50!black!50, inline]{#1}}
\newcommand{\todoResearch}[1]{\todo[color=red!50!black!20,bordercolor=red!50!black!50, inline]{#1}}
\newcommand{\todoImage}[1]{\todo[color=orange!50!black!20,bordercolor=orange!50!black!50, inline]{#1}}
\newcommand{\todoQuestion}[1]{\todo[color=blue!50!black!20,bordercolor=orange!50!black!50, inline]{#1}}

%----------------- FARBEN DEFINIEREN ----------------- %
\definecolor{gray}{gray}{0.95} % Listingsbackground

%----------------- LAYOUT SETZEN ----------------- %
%3, 2, 2.5, 2
\geometry{left=3cm, right=3cm, top=3cm, bottom=3cm}
\linespread {1.25}\selectfont %1.25 da er von Haus aus 1.2 ist und 1,25 * 1,2 = 1,5 isch


%-------##-------##-------##------- ANFANG INHALT -------##-------##-------##-------%
\begin{document}


\pagenumbering{Roman} % Seitennummer

%----------------- DECKBLATT -----------------%
%----------------- KONFIGURATION ----------------- %
\pagestyle{empty} % enthalten keinerlei Kopf oder Fuß 

%----------------- HDA FBI Logo ----------------- %
\begin{figure}[t]
	\centering
	\includegraphics[width=0.6\textwidth]{Abb/logo_fbi}
\end{figure}

%----------------- INHALT ----------------- %

\begin{center}
\Large Hochschule Darmstadt \\
\normalsize \textsc{- Fachbereich Informatik -} \\

% Whitespace
\vspace{85 pt}

\Huge Analyse sicherheitsrelevanter Informationen in industriellen Netzwerken für den Einsatz eines SIEM Systems \\ 
\normalsize
\vspace{20 pt}

Abschlussarbeit zur Erlangung des akademischen Grades \\ 
Master of Science (M.Sc.) 

\vspace{55 pt}


vorgelegt von \\
\vspace{5 pt}
Niklas Breuer\\
727905
\vspace{85 pt}

\begin{tabular}[h]{p{4cm}l l}
	Referent: & Prof. Dr. Oliver Weissmann\\
	Korreferent: & Jürgen Zorenc
\end{tabular}


\end{center}
 
%-----------------ERKLAERUNG---------------%
%----------------- KONFIGURATION ----------------- %
\pagestyle{empty} % enthalten keinerlei Kopf oder Fuß

\chapter*{Erkl\"arung} % (fold)
\label{cha:Erkl\"arung}

\normalsize
Ich versichere hiermit, dass ich die vorliegende Arbeit selbstst\"andig verfasst und keine anderen als die im Literaturverzeichnis angegebenen Quellen benutzt habe. 
Alle Stellen, die w\"ortlich oder sinngem\"a\ss \ aus ver\"offentlichten oder noch nicht ver\"offentlichten Quellen entnommen sind, sind als solche kenntlich gemacht.
Die Zeichnungen oder Abbildungen in dieser Arbeit sind von mir selbst erstellt worden oder mit einem entsprechenden Quellennachweis versehen.
Die Arbeit ist in gleicher oder \"ahnlicher Form noch bei keiner anderen Pr\"ufungsbeh\"orde eingereicht worden.
 
\vspace{5em}
 
\begin{flushright}
\begin{table}[ht]
\begin{tabularx}{\textwidth}{Xp{7cm}}
Darmstadt den \today & \tabularnewline	\cline{2-2} \addlinespace
& \centering{Niklas Breuer}
\end{tabularx}
\end{table}
\end{flushright}
 
%----------------- ABSTRACT -----------------%
%----------------- KONFIGURATION ----------------- %
\pagestyle{empty} % enthalten keinerlei Kopf oder Fuß


\chapter*{Abtract} % (fold)
\label{cha:abtract}
	
	\todoForm{Abstract formulieren}


 
 
\listoftodos 
 
%----------------- VERZEICHNISSE -----------------%
\tableofcontents % Inhaltverzeichnis

\pagestyle{plain} % zurueck setzen von roemische seitenanzahl
\pagenumbering{arabic}


%----------------- KAPITEL : EINFÜHRUNG  ----------------- %
\chapter{Einleitung}
\label{cha:Einleitung}
%Kontext 
\section{Kontext}
 
Ein grundlegendes Konzept für die Sicherung eines Netzwerkes ist das Wissen um den Zustand des Netzwerkes und seiner Elemente. In modernen Unternehmensnetzwerken mit tausenden Elementen, Schnittstellen und Abhängigkeiten ist es schwierig, jederzeit den Überblick zu behalten. Aus diesem Grund nutzen Unternehmen \textit{Security Information and Event Management (SIEM)} Systeme, die für die Analyse protokollierter Ereignisse und der Erkennung von potentiellen Bedrohungsszenarien verwendet werden.  

Die Industrie 4.0 ist ein Konzept, mit deren Hilfe die nächste Stufe der industriellen Revolution erreicht werden soll. Um dieses Ziel zu erreichen, sollen \textit{Smart Factories} erstellt werden, in denen Anlagen miteinander sowie Produkte mit den entsprechenden Anlagen kommunizieren können. Für dieses Vorhaben besteht ein steigendes Interesse an vertikaler und horizontaler Integration, d.h. der stärkeren Vernetzung der einzelnen Komponenten / Anlagen miteinander sowie der Anbindung industrieller Produktionsnetzwerke an das WAN (Internet). 
Dabei ist jedoch zu berücksichtigen, dass industrielle Produktionskomponenten, die auch mit dem Begriff \glqq cyber-physisches System\grqq{} bezeichnet werden, nicht für die Anbindung an externe Netzwerke entworfen wurden, sondern mit den Zielen Safety (Sicherheit der Mitarbeiter), Verlässlichkeit und Effizienz. Aus diesem Grund sind weder Komponenten noch die verwendeten Kommunikationstechnologien und -protokolle ausreichend gegen Bedrohungen aus dem Internet abgesichert. Dabei ist der Schutz dieser Systeme sogar besonders kritisch zu bewerten, da erfolgreiche digitale Angriffe sich aus der digitalen auf die physische Welt übertragen und materielle und finanzielle Schäden sowie eine Gefahr für die Beschäftigten der Anlagen darstellen können. 
Allerdings ist der Schutz dieser Systeme auf Grund mehrerer Faktoren problematisch. Die Zertifizierungsnotwendigkeit für die sichere Ausführung und der teils sehr alten Technologien und -standards ist nicht ohne weiteres möglich, da neben den Anforderungen der IT-Sicherheit für kritische Systeme vor allem auch die Anforderungen der Anlagentechnik gewahrt bleiben müssen. Die verwendeten speicherprogrammierbaren Steuerungen (SPS) sind für die Kontrolle der einzelnen Fertigungsprozessschritte auf eine sichere Funktionalität, jedoch nicht auf Manipulationssicherheit durch externe Quellen ausgerichtet. Basierend auf den Anforderungen innerhalb der industriellen Fertigungsnetzwerke sind Logik und Technologie auf Verlässlichkeit (d.h. Determinismus / Fehlerfreiheit), Reaktionsschnelligkeit und Robustheit gegen äußere physikalische Einflüsse der Anlagenumgebung ausgelegt. Diese Anforderungen und die ursprüngliche Entwicklung der Automation von logischen Steckkombinationen hin zu programmierbaren Steuerungen ergeben knappe Ressourcen im Sinne der Kommunikationsbandbreite, Speicher und Rechenleistung. Programmerweiterungen können durch zusätzliche Programmlogik Fehler und unberechenbares Verhalten des physikalischen Equipments mit sich bringen. Fehlende Standards ergeben zusätzlich eine Netzwerklandschaft, in der viele verschiedene herstellerspezifische Kommunikationsprotokolle, Geräte, Gerätetypen und Kommunikationstechnologien vertreten sind. Zwar werden Versuche und Anstrengungen unternommen Standards zu schaffen, allerdings ist der aktuelle Stand noch nicht soweit.

Diese Faktoren erschweren die Anwendung von typischen Sicherheitselementen wie den erwähnten SIEM Systemen ungemein. Neben der Notwendigkeit der Anpassung an die Anforderungen der Fertigungssysteme behindern die knappen Ressourcen und Limitierungen der Erweiterbarkeit und Aktualisierung der Anlagen die Anwendung der etablierten Sicherheitskonzepte aus den Unternehmensnetzwerken.

\section{Zielsetzung und Zweck}
In dieser Ausarbeitung soll diese Problematik im spezifischen Fall der SIEM Systeme untersucht werden. SIEM Systeme benötigen Informationen der Netzwerkelemente (sowohl der Endgeräte als auch der Netzwerkgeräte und der Kommunikationsmechanismen- und Schnittstellen). Ist die Verfügbarkeit der Informationen limitiert, kann das SIEM System auch nur im Rahmen der verfügbaren Informationen ungewöhnliches Verhalten der Systeme und Bedrohungsszenarien erkennen und die Sicherheitsoperatoren bei ihrer Arbeit unterstützen. Für die Absicherung eines industriellen Netzwerkes im Kontext der Industrie 4.0 ist eine Analyse der verfügbaren Informationsquellen notwendig. Basierend auf der bereits erläuterten Problematik wird jedoch die Hypothese aufgestellt, dass die typischen Komponenten und Technologien eines industriellen Netzwerkes kritische Informationslücken aufweisen, welche die Überwachung und Analyse kritischer Aspekte verhindern. Dies ist die Kernthese dieser Arbeit, die es zu untersuchen gilt. Für die Beantwortung dieser These werden die folgenden Forschungsfragen untersucht:
\begin{itemize}
\item Welche Informationen lassen sich jeweils aus Unternehmens- und Industrienetzwerken gewinnen?
\item Welche Unterschiede bestehen bzgl. der verfügbaren Informationen im Vergleich zwischen den Unternehmenstypen?
\item Sind die gefundenen Unterschiede als kritisch für die Überwachung industrieller Netzwerkes einzustufen?
\end{itemize}

Der Fokus für die Untersuchung der Unterschiede zwischen den Netzwerken soll auf das Fehlen von Informationskategorien in industriellen Netzwerken gesetzt werden. Die Elemente dieser konkreten potentiellen Untermenge der Unterschiede werden im Rahmen dieser Ausarbeitung als \textit{Informationslücken} bezeichnet.

\section{Vorgehen}
Als Startpunkt für die Suche nach den potentiellen Informationslücken in industriellen Netzwerken wird eine Referenzmenge benötigt. Diese Referenzmenge ergibt sich für SIEM Systeme aus bereits etablierten Implementierungen in Unternehmensnetzwerken. Aus diesem Grund ist der erste Schritt die Analyse der Informationsmenge eines Unternehmensnetzwerkes. Zu diesem Zweck werden zunächst Grundkenntnisse über das SIEM Konzept und dessen Anwendung sowie über Elemente und Technologien in Unternehmensnetzwerken und industriellen Netzwerken zur Verfügung gestellt. 

Die erste Herausforderung bei der folgenden Analyse ist die Eingrenzung der Informationsmenge. Durch eine sehr große Anzahl an unterschiedlichen Systemen, Protokollen, Schnittstellen und herstellerspezifischen Zusatzeigenschaften muss die Analyse auf ein handhabbares Modell beschränkt werden. Für die Analyse werden ein repräsentatives Modell definiert und die einzelnen Elemente genauer betrachtet. Ziel der Analyse ist es zum einen, die Informationsquellen, etwa in Form von Protokolldateien, zu identifizieren und diese in Kategorien auf Basis ihres Inhaltes einzuordnen. Dies ermöglicht die Eingrenzung und Fassung der Informationsmenge. 

Der gleiche Ansatz wird für die Analyse des industriellen Netzwerkes gewählt. Im Unterschied zu der vorherigen Analyse wird jedoch der Versuch unternommen, die Informationen den zuvor definierten Kategorien des Unternehmensnetzwerkes zuzuordnen. Basierend auf dieser Zuordnung entsteht eine notwendige Vergleichsbasis, auf der die beiden Netzwerkarten miteinander verglichen werden können.

Die Ergebnisse des Vergleiches müssen jedoch nicht zwangsläufig auf eine kritische Informationslücke schließen lassen. Durch die unterschiedlichen Zwecke und Fähigkeiten der Netzwerkelemente und Kommunikationstechnologien können Unterschiede in der Informationsmenge auch bspw. auf nicht unterstützte Funktionalitäten oder Dienste hinweisen. Um die Bedeutung von Unterschieden bzgl. der Notwendigkeit für die Überwachung zu ermitteln, müssen diese Unterschiede diesbzgl. bewertet werden. Die Bewertung wird anhand des populären Beispiels \textit{Stuxnet} für einen Angriff auf eine industrielle Anlage durchgeführt. Die Analyse der Aktionen des Stuxnet Angriffes erlaubt Rückschlüsse darauf, welche Informationen eines industriellen Netzwerkes notwendig sind, um Angriffe wie diesen zu erkennen. 
Ein Abgleich der gefunden Unterschiede mit den Analysepunkten des Stuxnetangriffes ergibt die Bedeutung der gefundenen Unterschiede.

Schlussendlich wird auf Basis dieser Erkenntnisse versucht, einen Lösungsansatz für eine konkrete Problemstelle unter bestimmten Annahmen zu formulieren.


\section{Struktur}
Die Struktur der Arbeit orientiert sich direkt an der beschriebenen Vorgehensweise. Kapitel 2 liefert die notwendigen Grundkenntnisse, die für das Verständnis der Untersuchung als notwendig erachtet werden. In den darauf folgenden Kapiteln werden die Analyse der Modelle des Unternehmensnetzwerkes (Kapitel 3) und des industriellen Netzwerkes (Kapitel 4) beschrieben. Die Vergleichsmethodik und die Ergebnisse des Vergleichs finden sich in Kapitel 5. Kapitel 6 beschreibt die Bewertung der Vergleichsergebnisse. Kapitel 7 beschreibt einen Lösungsansatz für eine ausgewählte Informationslücke.

%----------------- KAPITEL : GRUNDLAGEN  ----------------- %	
\chapter{Grundlagen}
\label{cha:grundlagen}


\section{IT Security Schutzziele}
\label{cha:IT Security Schutzziele}

In der IT-Sicherheit werden alle Ziele bezeichnet, die für die Sicherung von Systemen und Kommunikation zwischen Systemen sichergestellt werden müssen.
\todoForm{Mehr Kontext der Schutzziele einfügen}

Dabei werden als Grundlage die folgenden Schutzziele betrachtet:
\begin{itemize}
\item Integrität: Bei der Kommunikation zwischen Systemen und der Speicherung bzw. dem Abruf von Daten auf einem System ist es wichtig, dass darauf vertraut werden kann, dass diese Daten nicht ohne Authorisierung verändert wurden. Daher befasst sich dieses Schutzziel mit der Korrektheit von Daten bzw. der korrekten Funktion eines Systems. Bei der Sicherstellung dieses Zieles geht es also darum die Veränderung von übertragenen oder gesicherten Daten und Software sicherzustellen.
\item Vertraulichkeit: Innerhalb einer Infrastruktur oder auf einem System existieren verschiedene Arten von Daten, darunter auch sensible Daten,  die von nicht-authorisierten Personen missbraucht werden können. Daher ist es notwendig den Zugriff zu Daten zu limitieren und dadurch diese Daten vor unauthorisiertem Zugriff zu schützen. Aktivitäten dieser Art werden dem Schutzziel Vertraulichkeit zugeordnet
\item Verfügbarkeit: Diesem Schutzziel werden alle Aktivitäten zugeordnet, die den Betrieb von Systemen und die Erhaltung von Kommunikationswegen sicherstellen. Darunter fallen z.B. Maßnahmen, die den Zugriff von Kunden und/oder Mitarbeitern auf einen Webservice sicherstellen, auch im Falle eines Angriffes, der es zum Ziel hat die Erreichbarkeit des Services zu unterbinden.
\end{itemize}
\todoForm{Schutzziele: Beispiel für Schutzziele formulieren}
Neben diesen Schutzzielen existieren auch noch weitere Schutzziele wie z.B. Authentizität, Nichtabstreitbarkeit, Zurechenbarkeit und Privatsphäre.


\section{Security Information and Event Management}
\label{cha:Security Information and Event Management}

\subsection{SIEM Konzept und Zweck}
\label{cha::SIEM Konzept und Zweck}
Im Umgang mit Security Information and Event Management (SIEM) Systemen werden oft die Begriffe „Information“, „Ereignis“ und „Daten“ verwendet. Um eine Grundlage für die Verwendung dieser Begriff in der vorliegenden Thesis zu geben, werden im folgenden diese Begriffe wie folgt definiert:
\begin{itemize}
\item Ereignis (Event): Unter einem Ereignis ist in diesem Zusammenhang das Auftreten von systemrelevanten Aktionen zu verstehen. Dabei kann zwischen Systemereignissen, z.B. das Laden eines Programmes, und Benutzerereignissen, z.B. Anschläge auf der Tastatur, unterschieden werden. 
\item Daten: Daten sind die Bausteine aus denen Ereignisse zusammengesetzt werden. Daten sind z.B. der Name des geladenen Programms oder der Wert einer Benutzereingabe.
\item Information: Eine Information ist in diesem Kontext die Interpretation verschiedener Ereignisse, die eine Aussage über den Zustand eines Systems ermöglicht.
\end{itemize}


Die IT Infrastruktur von Unternehmen umfasst eine große Menge an Elementen wie Server, Arbeitsstationen, Netzwerkgeräte, Sicherheitssysteme und mobile Endgeräte (z.B. Laptops und Mobilfunkgeräte). Um diese Infrastruktur zu schützen reicht es nicht diese hinter einem Schutzwall zu positionieren, es ist auch wichtig zu erfassen welche Aktionen von wem wie und von wo innerhalb des Netzwerkes ausgeführt werden. Daher ist es wichtig, dass es Informationsquellen gibt aus denen diese Informationen ausgelesen und bewertet werden können. In Unternehmens-netzwerken wird dies typischerweise von den Elementen selbst durchgeführt. So befinden sich z.B. auf einem Server mit den Betriebssystem Windows verschiedene Log-Dateien, die Informationen darüber speichern welcher Anwender sich zu welchem Zeitpunkt angemeldet hat, wie viele Versuche für die Eingabe des Passwortes verwendet werden und welche Prozesse von diesem Anwender ausgeführt wurden. 
Die Analyse dieser Log-Dateien kann fachkundigen Administratoren Aufschluss darüber geben welche Aktionen von dem Benutzer oder dem System durchgeführt wurden. Eine Herausforderung bei dieser Analyse ist die Bewältigung der schieren Menge an Informationen und das Filtern relevanter Ereignisse. Selbst eine geringe Anzahl an Systemen kann eine große Menge an Daten produzieren welche von einem Administratorenteam ohne Hilfe von Werkzeugen nicht zu bewältigen sind. Mit Hilfe technischer Werkzeuge können die Daten in Log-Dateien gefiltert werden. Die Suche nach sicherheitsrelevanten Informationen kann dabei automatisiert werden und Administratoren können über Anomalien im Verhalten der Systeme informiert werden. 

Allerdings bietet diese Vorgehensweise auch Nachteile. So ist die Definition der Regeln nach denen Werkzeuge Log-Dateien durchsuchen eine komplexe Aufgabe, denn diese Regeln müssen eng genug gefasst sein um die Anzahl der Sicherheitsmeldungen verwaltbar zu halten, aber auch weit genug um potentiell bedrohliche Situationen zu erkennen. Um eine bedrohliche Situation bewerten zu können müssen Administratoren nicht nur in der Lage sein die Meldungen des Werkzeuges zu verstehen, sondern diese auch im korrekten Kontext einordnen zu können. Diese Einordnung kann in komplexen Netzwerken sehr schwierig sein.

Um diese Einordnung zu vereinfachen und die Anzahl an Falschmeldungen („False Positives“) zu reduzieren ist es also notwendig nicht nur ein Element sondern verschiedene Elemente im Zusammenspiel zu betrachten und die Informationen dieser Netzwerkelemente zu verknüpfen, also einen Zusammenhang von Informationen von verschiedenen Elementen zu erfassen. Für die Unterstützung der Administratoren bei dieser Aufgabe wurden Security Information and Event Management (SIEM) Systeme entwickelt.


%Beschreibung des Security Information and Event Management (SIEM) Konzeptes
Ein SIEM System ist also ein System, welches Informationen aus verschiedenen Quellen extrahiert, die einzelnen Informationsquellen analysiert und die gewonnenen Informationen in einen Zusammenhang bringt. Durch Aggregation von ähnlichen Daten und Korrelation dieser Daten lassen sich einzelne Alarmmeldung in einen Zusammenhang setzen und neue Informationen über den Wert und den Kontext der Meldungen gewinnen. Dadurch ist es möglich den Zustand eines Netzwerkes und seiner Elemente über ein zentrales System zu überwachen und kritische Situationen zu erkennen, die durch das Betrachten einzelner Elemente nicht erkennbar sind. Das SIEM System bildet damit eine zentrale Verwaltungsschicht oberhalb der Netzwerkelemente.

Um diese Funktion auszuführen sind verschiedene Schritte notwendig. Die Aufgaben teilen sich dabei in zwei Bereiche auf: Security Information Management (SIM) und Security Event Management (SEM). SIM ist eine Unterkategorie des Log Management Feldes, das heißt es umfasst Erfassung, Extraktion, Transfer und Sicherung von sicherheitsrelevanten Informationen in Log-Dateien. SEM umfasst Funktionen für die Analyse und Verknüpfung von Ereignissen in Echtzeit.

Der erste Schritt ist die Extraktion der Daten von den Netzwerkelementen. Dies kann entweder durch Anbindung einer geeigneten Schnittstelle durchgeführt werden (z.B. Meldungen von Sicherheitssystemen wie einer Firewall oder einem Intrusion Detection System (IDS)) oder durch Extraktion der Informationen über Kollektoren in Form von Software Agents. Der zweite Schritt betrifft die Übertragung der extrahierten Informationen. Diese muss sicherstellen, dass die Informationen vollständig und unverändert übertragen werden, da veränderte Informationen die Informationsgrundlage verändern auf der das SIEM System seine Analyse durchführt. 

Da nicht jedes System gleich ist und auch Log-Dateien und -Formate sich deutlich unterscheiden können, müssen die vom SIEM empfangenen Informationen vorverarbeitet und in ein einheitliches Format umgewandelt werden. Dieser Schritt wird als „Normalisierung“ bezeichnet. Die normalisierten Daten können dann zentral gespeichert und für statistische Langzeitanalysen und statistische Verwertungen gespeichert werden. Dabei ist es wichtig darauf zu achten, dass die abgelegten Daten nicht verändert werden können und das Änderungen nachvollziehbar sind. Da die Datenmenge abhängig von der Größe und Komplexität sehr groß sein kann und die Ressourcen des SIEM Systems für die Analyse und Bewertung begrenzt sind, werden ähnliche Daten in einem Zwischenschritt zusammengefasst, sodass die zu analysierende Menge an Daten deutlich reduziert wird. Man spricht dabei von der Aggregation der Daten. Die Aggregation kann z.B. dadurch erfolgen, dass Ereignisse desselben Ursprungs und desselben Inhalts als eine Meldung zusammengefasst und mit einem Zähler versehen werden, der die Anzahl der Meldungen wiederspiegelt.
Im nächsten Schritt werden die aggregierten Daten dann durch ein Regelwerk analysiert. Dieses Regelwerk wird manchmal auch als „Rule Correlation Engine“ bezeichnet. Dabei werden verschiedene Meldungen mit Regeln abgeglichen. Wird eine Regel als erfüllt angesehen, wird eine Alarmmeldung an verantwortliche Administratoren ausgegeben, sodass weitere Maßnahmen ergriffen werden können.  Die Bewertung ob eine Regel erfüllt wurde hängt von der verwendeten Korrelationstechnik ab. So müssen z.B. eine Reihe von Bedingungen erfüllt werden, damit eine Regel als erfüllt angesehen wird. So könnte z.B. eine Regel beinhalten, dass ein Alarm ausgegeben wird, wenn mehrere nicht-erfolgreiche Anmeldungsversuche von einem inaktiven Account von einer nicht-registrierten IP-Adresse protokolliert wurden. Die Bildung solcher Regelwerke kann sehr komplex sein und durch verschiedene Techniken gebildet werden.


%Grenzen
In der Praxis werden SIEM Systeme u.a. in Security Operations Center (SOC) eingesetzt. Das SIEM hilft in diesem Kontext nicht nur den Administratoren für die Verwaltung, sondern ersetzt auch die Verwaltung und Analyse über verschiedenen UIs unterschiedlicher Sicherheitsprodukte (z.B. von Firewalls, IDSs, Anti-Virus Software). 
Hier zeigen sich neben den Vorteilen von SIEM Systemen auch aktuelle Grenzen. So kann ein SIEM nur auf der Informationsbasis operieren, die in der jeweiligen Umgebung zur Verfügung steht. So kann sich das Fehlen von Informationen zum Kontext der Meldungen auf die Auswertung auswirken, da fehlende Informationen zu ungenauen Analysen oder False Positives bzw. False Negatives (kritische Meldungen, die als harmlos klassifiziert werden) führen können. Eine andere Limitierung sind Abfragezeiten von gespeicherten Langzeitdaten für die Echtzeitanalyse durch die entweder die Analyse oder die Menge der Daten begrenzt wird. Auf der technischen Seite ergeben sich zudem Herausforderungen entlang der Funktionskette eines SIEM Systems bei der Datenextraktion aus proprietären Formaten, der effektiven und sicheren Speicherung der Daten sowie der Analyse und Erstellung von Korrelationsregeln in komplexen Systemen. 


\subsection{SIM \/ Log Management}
%Kontext
Security Information Management (SIM) bezeichnet die Verwaltung von Log-Dateien im SIEM Kontext und bildet damit eine Untermenge des Log Management. Sicherheitsrelevant sind in diesem Falle alle Log-Dateien, die Auskunft über den Zustand eines Systems oder einer Kommunikation zwischen Netzwerkelementen geben können. Dazu zählen sowohl Log-Dateien von Betriebssystemen und Programmen als auch von Netzwerkswitches und anderen Elementen, die den Datenverkehr zwischen Teilnehmern im Netzwerk aufzeichnen. Eine Log-Datei besteht aus Ereignissen, die Daten zu Aktionen (u.a. Zeitpunkt, ausführender Benutzer, ausführender Prozess (ID), …) enthalten. Diese Daten können interpretiert werden und liefern Informationen über den Zustand des Systems zum bestimmten Zeitpunkt.

%Definition
Log Management umfasst die Sammlung, Übertragung, Normalisierung, Zentralisierung und Aufbewahrung der Menge an Log-Dateien. In großen Unternehmen kann diese Menge mehrere hundert Gigabyte umfassen. Das Verwalten der Log-Dateien wird aus mehreren Gründen in Unternehmen als außerordentlich wichtig angesehen. Zum einen beinhalten Log-Dateien Informationen über den Zustand der Umgebung und können damit Rückschlüsse auf potentiell schädliche Aktionen zulassen, zum anderen können durch diese Informationen auch die Einhaltung von geltenden Richtlinien und Anforderungen gegenüber Kunden und Institutionen belegt werden. Daraus folgt, dass sichergestellt werden muss das die Informationen nicht nur korrekt aus den verschiedenen Quellen extrahiert werden können, sondern auch die Umwandlung und ggf. Veränderung der Daten dokumentiert werden muss um die Integrität der Informationen zu gewährleisten. Dies spielt u.a. für die Sicherheitsanalyse eine wichtige Rolle, da nur mit den richtigen Informationen die korrekten Rückschlüsse auf die Sicherheit der Umgebung geschlossen werden können. So lassen sich u.a. mit forensische Analysen rückwirkend Aktionsketten, die zu einer Richtlinienverletzung oder eines nicht-authorisierten Eindringens in das System geführt haben, rekonstruieren. Darüber lassen sich aus den Informationen Grundlagen für den Normalzustand eines Systems ablesen und entsprechende Regeln formulieren.
\todoForm{Log Management: Unterschied zwischen Log Management Systemen und SIEM Systemen herausstellen/formulieren}

Log-Dateien können aus fast jedem System gewonnen werden. Dazu zählen u.a. 
\begin{itemize}
\item Anti-Malware und Anti-Virus Systeme
\item Intrusion Detection Systems / Intrusion Prevention Systems
\item Remote Access Software
\item Web Proxieserver
\item Vulnerability Management Software 
\item Network Access Control (NAC) Server
\item Firewalls
\item Router
\end{itemize}

Eine interessantes Thema in diesem Zusammenhang ist die Frage welche Log-Dateien und Event-Daten sicherheitsrelevant sind. Moderne und komplexe IT-Umgebungen produzieren mehrere hundert Gigabyte an Event-Daten auf einer täglichen Basis. Daher ist die Definition sicherheitskritischer Elemente (Systeme, Kommunikationspfade, Applikationen, ...) eine wichtige Aufgabe.
Dazu können unter anderem diese Elemente gehören:
\begin{itemize}
\item Logs der Security Controls (z.B. Firewalls, Intrusion Detection System\/Intrusion Prevention System, Data Loss Protection)
\item Logs der Netzwerk Infrastruktur (z.B. Domain Name Service (DNS) Server, Dynamic Host Configuration Protocol (DHCP) Server, VPN Logs)
\item Informationen über die Infrastruktur aus anderen Quellen (z.B. Informationen zu Systembestand und Netzwerksegment eines Elements)
\item Informationen über das Unternehmen
\end{itemize}

\todoQuestion{Log Management: Log Management Architektur einfügen?}
\todoForm{Log Management: Beispiel für Log Management Architektur einfügen}

%Herausforderungen
Eine große Herausforderung des Log Management ist das Schaffen einer Balance zwischen limitierten, verfügbaren Ressourcen für das Verwalten der Log-Dateien und der Menge an zu verarbeitenden Log-Dateien pro Zeiteinheit. Dabei spielen nicht nur die potentiell große Menge an Dateien in einer komplexen Umgebung eine Rolle, sondern auch der Umgang mit Inkonsistenzen zwischen vergleichbaren Logs (etwa in Bezug auf den Zeitstempel), der Umwandlung aus verschiedenen Formaten und das Wachstum an Daten mit dem Hinzufügen von weiteren Systemen in die Umgebung.

\subsection{Kollektoren}
%Kontext
Die Extraktion, oder auch Sammlung, der Daten von den Elementen des Netzwerkes bildet die Grundlage auf der das SIEM System (und ein Log Management System) funktioniert. Die Daten werden von vielen verschiedenen Geräten extrahiert, die eine gewisse gemeinsame Grundmenge an Daten bieten, darüber hinaus aber auch weiteren Kontext abhängig von Applikation, Betriebssystem oder Gerät selbst. Daher muss eine Schnittstelle geschaffen werden, die diese proprietären Datenformate auslesen und zu einem einheitlichen Format umwandeln kann. Dies ist die Aufgabe der Kollektoren.
\todoForm{Kollektoren: Absatz zu allgemeinem Extraktionsprozss (Source -> Log File -> Collector?)}

%Definition
Ein Kollektor ist demnach ein Service, oder auch Software Agent, der diese Aufgabe übernimmt. Dabei sind verschiedene Elemente zu betrachten, wie die verfügbaren Daten, Ressourcen der Datenquelle, Kommunikation zwischen Datenquelle, Kollektor und SIEM System sowie der Ansatz für einen bestimmten Gerätetyp. So stellt etwa das Protokoll, das für die Kommunikation mit der Datenquelle genutzt wird, eine Art Sprache dar, in der die Kommunikationspartner miteinander kommunizieren. Diese Kommunikationsform muss vorher vereinbart worden sein und dabei kann auf eine große Variation an Kommunikationsprotokollen zurückgegriffen werden. Der Kollektor für den entsprechenden Geräte-, Betriebssystem- oder Applikationstyp muss natürlich dieses Protokoll unterstützen. Alternativ dazu können Daten auch in weit verbreiteten Formaten wie etwa im „syslog“ Format abgelegt und durch den Kollektor oder eine Schnittstelle übermittelt werden.
Kommunikationsmodelle
Bei dem Ansatz der Datensammlung kann bzgl. SIEM Kollektoren zwischen dem Agent-basierten und dem Agent-losen Ansatz unterschieden werden. Der Agent-basierte Ansatz wird durch die Installation der Kollektoren auf den jeweiligen Netzwerkelementen aufgebaut. Dadurch ergibt sich eine dezentrale Verarbeitung, die zum einen Ressourcen des SIEM System spart, zum anderen aber auch einen erhöhten Verwaltungsaufwand für Kollektoren auf den verschiedenen Elementen des Netzwerkes bedeutet. Der Agent-lose Ansatz hingegen nutzt Kollektoren als Schnittstelle, die vom zentralen SIEM System aus mit den Komponenten über Schnittstellen kommuniziert. Hier wird die Verwaltung der Kollektoren erleichtert, aber es muss auch die Netzwerklast berücksichtigt werden, die durch die Kommunikation zwischen den SIEM Kollektoren und den Datenquellen erzeugt wird. Zudem muss über die entsprechenden Schnittstellen sichergestellt werden, dass Daten nicht auf der Datenquelle selbst oder während der Übertragung verändert werden.

Orientiert an den Aussagen aus der Industrie wird mehr und mehr zum dezentralen Ansatz tendiert, da die Agenten zusätzliche Funktionalitäten bieten können sowie Herausforderungen im Bereich der Regelung von Remote-Administrations Rechten, Security Compliance, Ressourcenmanagement und der Speicherung von hoch privilegierten Zugriffsdaten zu verschiedenen Systemen auf einer zentralen Instanz. Allerdings ist der Einsatz von Agenten auf Elemente mit einem entsprechenden Betriebssystem beschränkt. Hardware-Elemente mit proprietären Betriebssystemen wie etwa Netzwerk-Geräte unterstützen den Agenten-basierten Ansatz nicht und müssen daher selbst die Informationen in einem für das SIEM verständlichen Format zu der SIEM Instanz schicken oder eine Schnittstelle für den administrativen Remote-Zugriff bieten.

\todoForm{Kollektoren: Zeige die Datenextraktion anhand eines Beispiels}
<Quelle: Alien Vault>
Key problem: Normalizing log data
"Breaking down those logs from many sources into their components, or normalizing them, 
is what allows the SIEM to search across logs from multiple devices and to correlate events 
between them"
--> D.h. das SIEM muss in der Lage sein Security Incidents aus der Korrelation der Informationsevent zu erkennen
--> Die Kernfrage ist: Ist das möglich? In welchem Grade?
Interessanter Punkt: "Asset discovery" --> Informationen dafür sind?

\subsubsection{Herausforderungen auf Basis von Log Dateien}
\todoForm{Kollektoren: Formulierung der Notizen zu den Herausforderung der Informationsextraktion aus Log Dateien}
- Aussage zu Log Correlation: "This is mainly caused by a lack of standards that can encapsulate all events in a coherent way."
				"As a result, correlating between logs produced by different systems that use different log formats has been difficult and infeasible in many cases"
- Weitere Erwähnung eines einheitlichen Formats (auch im obrigen Paper): CEE (Common Event Expression) standard
- Aussage zu Event Sources: "A network-enabled device can further be subdivided into hosts and other low-level hardware"
	--> hosts into OS and Software Level
- Paper beschreibt die Art der Informationen in Logs von den vorher genannten Typen
	--> OS: "This includes, but is not limited to, the hardware configuration and state, allocated system resources and many types of security-related operations, such as authentication and authorization"
	--> Software: Conditions of the software
	--> Hardware (includes network infrastructure devices (switches, routers) and peripherals (VOIP telephones, printer))
		- "Event information from infrastructure devices are especially valuable for SIEM/IDS systems, because they control access to nodes in the network"
- Aussage zu "Event Persistence": "The employed mechanisms are usually the persistence of events directly into files, databases or delegation of persistence to a remote server (e.g. a syslog server)"
	--> "Log files provide a challenge in the sense that, often, no real structure is provided because there is typically only a single stream of characters"
	--> Serialisation of multiple events is complex
		 
	--> Differences in the separation technique of multiple event logs in a single file
		--> Line-by-Line Separation
		--> Delimiter Separation
		--> Pattern Separation (single and multiple lines, every new event records starts with a timestamp)
		--> Structuring (e.g. by XML objects (Microsoft event log format)
	--> Event Record Serialisation
		--> Structured log format --> Every Information has a clear location and meaning; machine-readable (Beispiele: CEE und IDMEEF)
		--> Unstructured log format --> No fields, instead written like language, potentially not machine-readable (Beispiel: MSG field of Syslog)
		--> Semi-Structured
	--> Event representation: 
		--> Single-Line: "Often a single event is represented using multiple log lines; Cases: Read and interpreted by humans, single line is too large (e.g. over a longer period of time)
		--> Multi-Line: "

- Aussage zu "Challenges in Event Normalisation"
	--> Level of details (different amount of fields between formats which can potentially be filled)
	--> Missing Structure (e.g. Syslog and Snort, challenging because there is no direct mapping to fields)
	--> Completeness (not enough information, e.g. log produceser, source \& target of an activity)
- Aussage zu Log Formaten: "Having a look at the ariety of log formats useed by event sources, it is self evident that there is no clear standard for logging and no agreement on what information is necessary to make a log useful for log analysis"
	--> Existing SIEMs: Splunk, ArcSight, Prelude, RSA envision----------






\subsection{Event Verarbeitung}
\label{cha:Event Verarbeitung}
%Kontext
Im Zuge der Normalisierung von Log-Dateien spielt auch das Zusammenfassen von Informationen eine wichtige Rolle. Die Aggregation, d.h. die Zusammenfassung, von ähnlichen Events hat einige Vorteile. So kann u.a. die Menge der zu übertragenden Informationen deutlich reduzieren und damit die Netzwerklast reduziert werden. Auch erleichtert es die Analyse der eingehenden Daten und kann damit die benötigte Menge an technischen und zeitlichen Ressourcen reduzieren. Zu guter Letzt können durch die Reduzierung der Datenmenge diese Daten formatiert schneller und mit geringerem Speicherverbrauch gespeichert und abgerufen werden.

%Typen von Techniken (grob)
\todoResearch{Event Verarbeitung: Thema ausbauen - Check: Was fehlt hier noch? -> Recherche falls notwendig -> Formulieren}
Bei der technischen Umsetzung der Aggregation lassen sich verschiedenste Ansätze finden. 
\begin{itemize}
\item Strategisch: Strategisch/mathematisch basiert, Wissens-basiert, Hash-basiert und Mining-basiert
\item Kommunikation: Strukturiert und Unstrukturiert
\item Verarbeitung: Hierarchisch, Durchschnitt, Sketches, Digests, Deterministisch
\end{itemize} 

\subsection{Security Event Management}
\label{cha:Security Event Management}

Security Event Management befasst sich mit der Echtzeit-Analyse von normalisierten Events, der Korrelation dieser Events sowie der Benachrichtigung von Sicherheitsadministratoren und der Darstellung relevanter Informationen. Man kann in diesem Zusammenhang auch von kontext-bewusster Überwachung sprechen.
Bei der Echtzeit-Analyse werden die extrahierten, normalisierten Events auf ihre Bedeutung untersucht. Dabei werden die Events einzeln für sich betrachtet und basierend auf den Daten(feldern) Informationen gewonnen. Diese Informationen werden bei der Korrelation der Events verwendet um einen Zusammenhang herzustellen. Dabei werden u.a. zusätzliche Informationen hinzugezogen, z.B. die Quelle der Informationen und Informationen zu der entsprechenden Quelle.
Dies ist notwendig um die Herstellung falscher Zusammenhänge zu minimieren. Die Herstellung der Korrelation basiert auf den verwendeten Techniken. So können Regelwerke auf verschiedene Art und Weisen hergestellt werden. Die Spannweite der Korrelationsregeln reicht dabei von simplen Wertevergleichen über Regeln mit verschiedenen Bedingungen hin zu der Erzeugung komplexer Szenarien durch Machine Learning und Big Data Ansätze.

\todoResearch{SEM: Check fehlende Inhalte --> ggf. kurze Recherche -> Formulieren}
- Aussage zu Correlation: Typically rule based, either network/host based or time/frequency based
- Aussage zu SOC Ablauf: 
	--> SOC analysts look at alerts and escalata fishy alerts to a level 2 analysis
	--> Number of alerts is so high that SOCs cannot analyse all, so alerts with higher priority are analysed earlier
		--> This approach could miss important alerts
- Aussage des Papers: Machine Learning can support analysts from within the workflow to reduce number of false positives and providing risk scores


\subsection{Features \&  Gemeinsamkeiten von SIEM-System-Anbietern}
\todoForm{SIEM Anbieter: Formulierung notwendig}
- Vendors: Cisco, RSA, Arc Sight, IBM, Q1Labs, LogLogic, NetIQ, CA
- "Gemeinsamkeiten": SIEM als Log Management Applikation, (Korrelation von Events) 
- Unterschiede: Log Formate, Kollektoreneinsatz, Kontext-Analyse, Korrelationstechniken


\section{Infrastruktur Office}

Für den Vergleich zwischen Unternehmensnetzwerk und industriellem Netzwerk ist es notwendig beide Netzwerke zu kennen. Dieser Abschnitt stellt eine Fundament zur Verfügung auf dessen Basis die Analyse gestartet werden kann.

\subsection{Geräte}
Ein Unternehmensnetzwerk besteht aus vielen verschiedenen Elementen. Dazu gehören natürlich auch die physischen Elemente die sich grob in drei Kategorien einteilen lassen:
\begin{itemize}
\item Endgeräte (PCs, Laptop, Handy, ...)
\item Server
\item Netzwerkgeräte (Switches, Router, Firewalls, Sicherheitselemente (z.B. Firewalls))
\end{itemize}

Damit ein Unternehmen seine Aufgaben ausführen kann werden verschiedene Server-Elemente benötigt. Ein Server ist ein zentraler Computerelement, welches mit mehreren anderen Elemente verbunden werden kann und entsprechenden Zugriff erlaubt. Zudem verfügen Server oft über leistungsfähigere Hardware und Speicherressourcen als ein Endgerät. Server werden benutzt um zentrale Applikationen bereitzustellen und zu verwalten. Dazu zählen neben Applikationen und (zugehörigen) Datenbanken auch Benutzerverwaltungssysteme und Sicherungscluster (Cluster := eine Menge an Servern, die durch eine zusätzliche Virtualisierungsschicht zu einem Element verknüpft werden).

Netzwerkgeräte dienen der Verbindung und/oder Überwachung der Kommunikation zwischen den einzelnen Netzwerkelementen. Diese Geräte existieren in unterschiedlichen Stufen und Fertigungen und sind angepasst für die jeweilige Aufgabe und Position im Netzwerk. So werden bspw. verschiedene Firewalls in verschiedenen Positionen im Netzwerk unterschiedlich konfiguriert um bestimmte Verbindungen zuzulassen oder zu blockieren.

Endgeräte werden von Mitarbeitern und Gästen des Unternehmens verwendet als Zugangswerkzeug zum Unternehmensnetzwerk. Basierend auf den Richtlinien und Sicherheitsvorgaben sind diese Geräte deutlich eingeschränkt. Die Benutzer auf den Endgeräten werden üblicherweise in eine Domäne eingebunden um eine zentrale Benutzerverwaltung zu ermöglichen.

\todoForm{Enterprise Networks: End-User Geräte tiefer beschreiben (Notizen im Kommentar)}
%Laptops, Mobile Endgeräte, stationäre Computer
\todoForm{Enterprise Networks: Server tiefer beschreiben (Notizen im Kommentar)}
%Webserver, Datenbankserver, Applicationserver, Proxyserver
\todoForm{Enterprise Networks: Netzwerkgeräte tiefer beschreiben (Notizen im Kommentar)}
%Router, Switches, Firewalls, IDS, IPS

\subsection{Architekturen \&  Anforderungen}
Unternehmen können in verschiedener Art und Weise in einer Architektur strukturiert sein, die auf das entsprechende Firmenmodell angepasst sind. 
Ein grundlegender Aufbau kann z.B. so aussehen:

\todoImage{Bild für typische Enterprise Network Architecture einfügen}
Zunächst besteht eine Verbindung zum Internet, diese wird durch gesicherte Gateways ermöglicht. Damit allerdings keine direkte Verbindung von diesen Gateways in das interne Netzwerk gewährt wird, wird eine zusätzliche Netzwerkzone zwischen das Gateway und das interne Netzwerk geschaltet. Diese Zone wird de-militarisierte Zone (DMZ) genannt. Innerhalb der DMZ werden Server angebunden, die aus dem Internet erreichbar sein sollen. Dies umfasst u.a. Webserver, die einen oder mehrere Webservices beherbergen oder Sprungserver (z.B. für eine VPN-Verbindung). Diese Server werden im Bezug auf ihre Kommunikationsfähigkeit limitiert, sodass keine Netzwerkverbindung von diesen Servern in das interne Netzwerk hergestellt werden kann, nur in umgekehrter Richtung ist die Herstellung einer Verbindung möglich. Das interne Netzwerk widerum kann in viele verschiedene weitere Netzwerkzonen und Domänen unterteilt sein. Diese Segmentierung dient der Strukturierung des Netzwerkes und eröffnet zusätzlich die Möglichkeit weitere Sicherheitsschichten in die Umgebung einfließen zu lassen. Das interne Netzwerk beherbergt Endgeräte von Mitarbeitern, Server mit Firmendaten, Datenbankserver und Datensicherungscluster. Alle Elemente in einer solchen Umgebung sowie die Kommunikation zwischen diesen Elementen können von verschiedenen Sicherheitssystemen überwacht werden. 

%Ausbauen
\todoForm{Enterprise Networks Architecture: Ausbauen - Netzwerkzonen, Beschreibung der Verschachtelung der Netzwerke, Spezifische Positionen (IDF, MDF, NOC)}

\subsection{Kommunikation}
In Unternehmensnetzwerken kommunizieren die Elemente auf Basis der Ethernet-Technologie und dem Internet Protocol (IP). Der zugehörige TCP/IP Stack bestimmt die grundlagen der Kommunikation durch die Zuweisung einer eindeutigen IP Adresse. Die Zuweisung dieser IP Adresse kann fest zugeordnet werden (z.B. für Server), manuell erfolgen oder automatisch über das Dynamic Host Configuration Protocol (DHCP) erfolgen, welches einem neuen Netzwerkelemente automatisch eine freie IP Adresse aus einem Adressenpool zuweist. Die meisten Elemente in einem Netzwerk werden einer sogenannten Domäne zugeordnet. Eine Domäne ist eine Zusammenstellung verschiedener Netzwerkelemente, die eine zentrale Benutzerverwaltung für die zugehörigen Elemente erlaubt. 
\todoForm{Beschreibung DNS hinzufügen}
Basierend auf der IP-Adresse kann eine verbindungsorientierte Kommunikation per Transmission Control Protocol (TCP) oder eine verbindungslose Kommunikation per User Datagram Protocol (UDP) erfolgen.  Ist es eine Priorität, dass eine Nachricht vollständig und in richtiger Reihenfolge übertragen wird, wird z.B. TCP verwendet. Auf weiteren zusätzlichen Schichten können weitere Funktionalität wie etwa eine kryptografische Verschlüsselung erfolgen. Die verschiedenen Schichten sind u.a. im OSI Layer Model festgehalten.
\todoForm{Enterprise Kommunikation - VPN \& Remote Access beschreiben}
\todoQuestion{Enterprise Kommunikation - Fehlt hier sonst noch etwas?}

\subsection{Bekannte Angriffsvektoren}
\todoForm{Enterprise Angriffsvektoren: Allgemeine Bedrohungen formulieren}
\todoResearch{Ggf. eine Ebene tiefer gehen, Recherche falls notwendig und formulieren}
Prinzipiell:
- Der Mensch als Ziel
- DDoS
- Malware \& (Spear) Phishing
- Ransomware

\section{Infrastruktur industrielle Produktionsnetzwerke}
\subsection{Geräte}
In industriellen Netzwerken werden verschiedene Komponenten und Technologien auf den verschiedenen Ebenen der Prozesskontrollstruktur eingesetzt. Zu den Hauptkomponenten zählen:
\begin{itemize}
\item Speicherprogrammierbare Steuerung (SPS)
\item Verteiltes Kontrollsystem
\item Human-Machine Interface
\item Industrial Ethernet Switch
\item Computer für den verwaltenden Zugriff auf angeschlossene SPSen
\item Sensor \& Aktuator
\end{itemize}


%Definition SPS
Eine speicherprogrammierbare Steuerung (SPS) ist ein Kontrollelement, welches für die Kontrolle eines Prozessschrittes genutzt wird.
%Aus dem Paper
Auf der SPS befinden sich sowohl eine Firmware als auch ein Programm. Die Firmware kontrolliert die Steuerung der SPS an sich. Das Programm wird in einen programmierbaren Speicher gelade. Das auf die SPS geladenen Programm verarbeitet Informationen, welche von den Eingangsmodulen empfangen werden. An die Eingangsmodule sind üblicherweise ein bestimmter Sensortyp angeschlossen, welcher Informationen über einen bestimmten Aspekt des Prozessschrittes liefert. Basierend auf der Implementierung des Programms, dies kann z.B. eine logische Sequenz, eine Zeitschaltung oder eine arithmetische Operation sein, wird durch diese Informationen ggf. über ein Ausgangsmodul ein Aktor (z.B. ein Motor) aktiviert, der ein Element des Prozessschrittes kontrolliert und Änderungen anstößt oder durchführt. Ein Programm kann mehrere Operanden kombinieren und in Zusammenhang setzen, inklusive Input, Output, Arguments, Counter, Timer und Function Blocks.

%Definition
Ein weiteres Element eine industriellen Netzwerkes ist ein verteiltes Kontrollsystem. Ein verteiltes Kontrollsystem ist für die Koordinierung und Überwachung eines Netzwerkes aus speicherprogrammierbaren Steuerung zuständig und kann empfangene Daten u.a. an einen Leitstand weiterleiten und/oder an einem Human-Machine Interface (HMI) sichtbar machen. Das HMI dient den Mitarbeitern für die Kontrolle eines Prozesses und die Zustände der ausführenden Elemente.

Für die Kommunikation über die Industrial Ethernet Schnittstelle existieren Industrial Ethernet Switches. Diese leiten den Datenpakete wie gewöhnliche Switche an die adressierten Empfänger weiter, sind aber darüber hinaus auch besonders für die Anforderungen in Fertigungsanlagen angepasst.

Zusätzlich zu dem Leitstand und den HMIs kann auch eine Verbindung zu SPSen über eine besondere Schnittstelle hergestellt werden, die es erlaubt Daten aus einer SPS auszulesen und ein neues Programm zu laden. 

Sensoren und Aktoren sind schließlich die Werkzeuge des Produktionsprozesses. Sensoren sammeln Daten über den Zustand der Produktion, etwa über die Höhe einer Flüssigkeit, die Temperatur eines Produktionselementes oder die Rotation eines Motors. Die Aktoren werden verwendet um genaue Anpassungen für den Prozess basierend auf der Daten der Sensoren und der Arithmetik der SPS-Programme vorzunehmen.

\subsection{Architektur \&  Anforderungen}
%Prozesskontrolle
Die Architektur eines industriellen Netzwerkes ist darauf ausgelegt, einen Fertigungsprozess zu kontrollieren. Dementsprechend existiert eine hierarchische Struktur, welche die Kontrolle des Prozesses von der Aufgabenannahme bis hin zur Kontrolle der einzelnen Schritte durch SPSen, Aktoren und Sensoren. Zu diesem Zweck wird auf die Disziplin der Prozesskontrolle zurückgegriffen. Die Prozesskontrolle ist eine Ingenieursdisziplin, die sich auf Architekturen, Mechanismen und Algorithmen für die Aufrechterhaltung der Funktionalität eines Prozesses beschäftigt. Allgemein spricht man bei Elementen innerhalb der Prozesskontrolle von industriellen Kontrollsystemen. Diese Kontrollsysteme können in verschiedenen Architekturen abgebildet werden.

Eine mögliche Form ist die Supervisory Control and Data Acquisition (SCADA) Architektur. In dieser Architektur existieren prinzipiell vier verschiedene Ebenen (von Level 4 (Verwaltung) zu Level 0 (Sensoren und Aktoren):
\begin{itemize}
\item Level 4 repräsentiert die Schnittstelle zum Unternehmensnetzwerk. Systeme auf dieser Ebene sind für die Produktionsplanung verantwortlich.
\item Level 3 repräsenteirt die Produktionskontrolle über die Anlage. Systeme auf dieser Ebenen kontrollieren den Produktionsprozess indirekt durch die Überwachung von Zielen und Produkten.
\item Level 2 beinhaltet weitere Systeme für die Beaufsichtigung wie z.B. den Leitstand. Die Elemente werden genutzt um Aktionen auf Level 1 und Level 0 zu kntrollieren. Human Machine Interfaces (HMI) gehören zu diesen Elementen und stellen eine grafische Benutzeroberfläche für Mitarbeiter zur Verfügung. Außerdem beinhaltet dieses Level das verteilte Kontrollsystem.
\item Level 1 ist die Ebene mit direkter Kontrolle auf die Prozessschritte über die angeschlossenen SPSen.
\item Level 0 beinhaltet Sensoren und Aktoren, die für die Kontrolle und Anpassung der Prozessschritte verwendet werden
\end{itemize}

\subsection{Kommunikationstechnologien}
%Kontext
In industriellen Netzwerken kommen verschiedene Elemente und Technologien zum Einsatz. Da Verfügbarkeit garantiert werden muss, werden Elemente wie SPSen über Jahre hinweg eingesetzt, sodass sich eine homogene Technologielandschaft ergibt, die unterschiedliche Kommunikationstechnologien unterstützt. So können bspw. ältere SPSen nur per Feldbus angesprochen werden, während neuere SPSen oder emulierte SPSen auch durch eine industrielle Version der Ethernet-Technologie angesprochen werden können. Während die Feldbustechnologie sich in der Produktionswelt durchgesetzt hat, wurde das Ethernet in Unternehmensnetzwerken umgesetzt. Heutzutage wird jedoch versucht auch die industrielle Version der Ethernet-Technologie stärker in der Produktion zu etablieren. Die Verbindung dieser Technologien für die Umsetzung des Konzeptes der Industrie 4.0 wird auch als Virtual Automation Network (VAN) bezeichnet.

%Feldbus
Feldbusnetzwerke sind standardisierte, aber auch proprietäre Netzwerke. In den Standards IEC 61158 und IEC 61784 existieren zehn verschiedene Konzepte. Sieben dieser Konzepte stellen eine eigene Protokoll Suite zur Verfügung, die anderen drei Konzepte basieren auf Ethernet Funktionalität. Beispiele für Protokolle sind etwa PROFIBUS oder DeviceNet.
Die Feldbuskommunikation stellt einen gemeinsamen Datenbus für angeschlossene SPSen bereit. Dabei können verschiedene Methoden der Kommunikationskontrolle angewendet werden um die Nutzung des gemeinsamen Kommunikationsbusses nutzbar zu machen. So wird etwa das sogenannte Master-Slave Prinzip angewendet, bei dem eine Master SPS den Kommunikationsfluss der anderen SPSen kontrolliert und Informationen von einer bestimmten SPS anfragt, die damit entsprechend die Erlaubnis zum Senden von Informationen erhält. Eine andere Möglichkeit ist das Master-Master Prinzip, bei dem ein Token von SPS zu SPS übertragen wird. Die SPS im Besitz des Tokens besitzt auch die Sendeerlaubnis.

%Industrial Ethernet
Eine andere Kommunikationstechnik wird in Form einer industriellen Version der Ethernet-Technologie verwendet, auch "Real-Time Ethernet" genannt. Dabei werden verschiedenen Ansätze verfolgt:
\begin{itemize}
\item Local soft real-time (TCP/IP Mechanismen mit ModbusTCP oder PROFINET CBA)
\item Deterministic real-time 
\item Isochronous real-time
\end{itemize}

Jeder der Ansätze kann unter bestimmten Bedingungen eingesetzt werden, basierend vorallem auf der Schnelligkeit der Übetragungen.

%Anforderungen
Diese Kommunikationstechnologien müssen bestimmte Anforderungen erfüllen. So muss eine sichere und robuste Kommunikation sowie schnelle Übertragungsgeschwindigkeiten garantiert werden, da eine verspätete Reaktion oder fehlerhafte Kommunikation eine Gefahr für die physischen Geräte und auch für die Sicherheit der Mitarbeiter darstellen kann. Daher ist die Verfügbarkeit das höchste Schutzziel, dass es zu erfüllen gibt.

\subsection{Kommunikationsprotokolle	}
\todoForm{Industrial Kommunikationsprotokolle: Formulierung nach Quellen des SPS Sec Paper}
\todoQuestion{Industrial Kommunikationsprotokolle: Habe ich hier alles an Materialien?}

\subsection{Bekannte Angriffsvektoren}
- Nicht viel bekannt
- Forscher haben verschiedene Elemente aufgedeckt
- Prinzipielles Prinzip

\section{Industrie 4.0 Konzept}
Das Konzept der Industrie 4.0 stellt einen Ansatz zur Steigerung der Flexibilität und Effektivität der Wertschöpfungsketten. Dazu sollen verschiedene industrielle Anlagen innerhalb einer Wertschöpfungskette über das Internet (WAN) miteinander verbunden werden. Dabei sollen sogenannte "Smart Factories" geschaffen werden, die sich besser auf die Wünsche der Kunden einstellen können. Wichtige Punkte dieses Konzeptes sind die horizontale und vertikale Integration. Zum aktuellen Zeitpunkt sind Unternehmens- und Industrienetzwerke voneinander getrennt und/oder durch eine Sicherheitsschicht voneinander getrennt, sodass keine Kommunikation zwischen Industrienetzwerk und Internet durchgeführt werden kann. Das Konzept der vertikalen Integration soll diesen Zustand verändern, sodass Produktionsverwaltung per Web Interface erfolgen kann. Gleichzeitig soll über das Konzept der horizontalen Integration eine stärkere Verknüpfung und Kommunikation der Elemente auf den verschiedenen Netzwerkebenen erreicht werden.


\subsection{Aktueller Sicherheitsstand von industriellen Produktionsnetzwerken (Schutzziele)}
Industrielle Netzwerke wurden ursprünglich als isolierte Umgebungen entwickelt, sodass Sicherheitsmaßnahmen sich stärker auf den physischen Zugang zu den Anlagen als auf die Sicherung der IT Systeme konzentrierten. Das höchste Schutzziel in diesen Umgebungen ist die Sicherstellung der dauerhaften Verfügbarkeit der Komponenten, d.h. das Garantieren der Robustheit und 

Durch das Auftreten von Stuxnet auf das iranische Atomprogramm und andere Angriffe auf industrielle Anlagen ist die Sicherung der Kommunikation und des virtuellen Zugriffes auf die Komponenten verstärkt als wichtiger Punkt anerkannt worden. Da in industriellen Netzwerken unterschiedliche Kommunikationstechniken existieren, die teilweise keinerlei Schutzmechanismen z.B. für die Verifizierung der Kommunikationspartner oder der Integritätsprüfung der ausgetauschten Informationen bieten, sind industrielle Netzwerke nach aktuellem Stand noch sehr anfällig für Angriffe. Dementsprechend ist die Sicherung dieser Netzwerke eine hohe Priorität, da bei einer Verbindung mit dem Internet ein potentieller Angriffsweg für Angreifer geschaffen wird, die bisher auf physischen Zugang oder das Einschleusen kompromitierter Speicher, wie z.B. USB Sticks, angewiesen sind. 
Neben den Schwachstellen der Kommunikationstechnologien ist die Sicherung bei gleichzeitiger Garantie der Verfügbarkeit eine weitere Herausforderung. So ist es nach aktuellem Stand nicht möglich, regelmäßige Sicherheitsupdates für SPSen aufzuspielen, da dies das Anhalten der Produktionsprozesse und damit signifikante Verzögerungen und finanzielle Verluste zur Folge hat. Darüber hinaus müssen Neuerungen an der SPS Firmware und/oder geladenen Programmen jeweils zertifiziert und überprüft werden um die Robustheit und Ausfallsicherheit zu garantieren. Dieser Zustand macht es schwierig gebräuchliche Sicherheitskonzepte aus Unternehmensnetzwerken in industriellen Netzwerken einzusetzen.

\subsection{Aktuelle Forschungsgebiete}
\todoResearch{Existierende Citavi Liste auswerten}
\todoResearch{Forschungsgebiete nochmal abdecken mit Suchlauf durch IEEEXplore, Elsevier, Springer und GooglePapers, Zeitrahmen: 2 Stunden}
\todoForm{Ergebnisse der Recherche Formulieren}

\begin{itemize}
\item Anomalieerkennung
\item Konzeptionen für sicherere Industrie 4.0 Komponenten
\item Untersuchung aktuell eingesetzter Komponenten und Kommunikationsprotokolle auf potentielle Sicherheitslücken
\end{itemize}

\section{Risiko Management}
\todoResearch{Risk Management: Gefundene Quellen bearbeiten}
\todoForm{Risk Management: Resultierende Notizen formulieren}
- Risk Management im Sicherheitskontext
- Critical Elements
- Abwägung: Risiko und Wahrscheinlichkeit	

%------------------KAPITEL : KOMPROMITTIERUNGSANSATZ ------%
\chapter{Analyse der Informationsquellen in einer typischen Unternehmensstruktur}
\label{cha:Analyse der Informationsquellen in einer typischen Unternehmensstruktur}

% Netzwerk so gut es geht nachbauen =)

\section{Beschreibung der Unternehmensarchitektur (Beispiel)}

%Kommentar: Vielleicht macht es Sinn erst die Elemente und dann die Architektur zu beschreiben...

%Ziel 1: Beschreibung der ausgewählten Kernfaktoren basierend auf der Grundlagenrecherche
%Ziel 2: Beschreibung der Architektur & Interaktionen als Grundlage der Analyse

%Architekturidee:
%Router <--> Firewall <--> DMZ (Webserver mit Webservice) <--> Firewall <--> Intranet (User PC) <--> Firewall <--> Restricted Zone (Database Server & SIEM)

/Das Ziel noch mehr herausstellen, was ist der Grundlegende Gedanke hinter der Erstellung des Models? 
//Einbinden, dass das Model auf der Basis von Netzwerkarchitekturen erstellt wurde
„Der Aufbau des Netzwerkes orientiert sich an typischen Elementen, die in einem Unternehmensnetzwerk zu finden sind. Dabei besteht die Notwendigkeit das Model in seiner Größe und Komplexität einzugrenzen um eine Analyse in sinnvollem Maße zu ermöglichen. Die Architektur und Einteilung in verschiedene Netzwerkzonen ist dabei von üblichen Sicherheitszonen abgeleitet. Während die Positionierung der Elemente keine besondere Rolle spielt im Bezug auf die ermittelbaren Informationen pro Element ist jedoch anzumerken, dass auch diese Informationen, z.B. in Form der Auswertung von Netzwerkadressen, einem SIEM-System nützliche Informationen zur Verfügung stellen können. Die Architektur des Models beginnt an seinem „Rand“, dem Zugang zum WAN (Internet) über einen Gateway-Router. Dieser Router ermöglicht die Weiterleitung von Netzwerkpaketen in und aus der angrenzenden Netzwerkzone, der de-militarisierten Zone (DMZ). In der DMZ sind für diese Model ein Apache Webserver auf der Basis des Betriebssystems „CentOS“ und ein E-Mail Server, bestehend aus dem Windows Webserver „Internet Information Service“ (IIS) und dem darüber liegenden E-Mail Server „Microsoft Exchange“. Diese Netzwerkelemente bieten die Möglichkeit Datenverkehr über die Protokolle HTTP und die Microsoft Schnittstelle MAPI (bestehend aus RPC & HTTP Datenverkehr) zu untersuchen, sowie Informationen aus Log-Dateien sowohl von einem Linux-basierten Server als auch von einem Windows-basierten Server zu analysieren. Alle Elemente werden durch „Managed Switches“ miteinander verbunden.
//Remote Access Server in Form von VPN, Protokolle? Traffic? Sinnvoll?

Die DMZ wird von der nächsten Zone, dem Extranet, durch eine Firewall überwacht, die eingehenden initiale Kommunikation blockiert. Firewalls gehören zu den Grundlagen der Sicherheit von Unternehmen und werden häufig platziert um Datenverkehr in und aus bestimmten Netzwerkzonen zu überwachen. Innerhalb der Netzwerkzone werden zwei Benutzer-PCs eingesetzt. Diese werden platziert um Abweichungen und andere Benutzerinteraktionen in die erhaltbare Informationsmenge zu integrieren. Die beiden Computer unterscheiden sich wie auch die Server der DMZ im Betriebssystem, sodass ein PC auf Basis von Windows und damit verbundener Benutzerverwaltung durch Active Directory enthalten ist, sowie ein PC mit einem Linux-basierten Betriebssystem. Zudem wird ein Netzwerkdrucker integriert.
//Gehören Printserver hier rein? Welche Zusatzinformationen bringt ein Printserver?

Als dritte Netzwerkzone wird als „Restricted Area“ bezeichnet. Innerhalb dieser Netzwerkzone werden verschiedene Windows-basierte Server eingesetzt um die Funktionalitäten eines Unternehmensnetzwerk zu simulieren. Die verschiedenen Servertypen wurden ausgewählt um erhaltbare Informationen aus unterschiedlichen, typisch genutzten Netzwerkprotokollen aufzuzeigen. Dazu gehören: 
\begin[itemize]
\item ein Microsoft SQL Datenbankserver 
\item ein Active Directory Domain Controller 
\item ein FTP-Server
\end[itemize]

Zusätzlich ist in dieser Zone als Ergänzung auch der SIEM-Server gesetzt.“

\subsection{Systeme}

%Frage: Warum lasse ich Elemente wie z.B. eine zentrale Benutzerverwaltung raus? Ist das nicht eigentlich relevant?
%Antwort: Die zentrale Frage hier ist: Welche Daten bekomme ich, im Bezug auf den normalen PC, aus einer Benutzerverwaltung wie z.B. einem Active Directory, die ich nicht aus einem Windowssystem bekomme?
%Um diese Frage zu beantworten müsste ich wissen nach welchen Informationen ich suche im Bezug auf die Benutzerverwaltung?
%Brainstorming: Login/Logoff von Benutzern (beide), Berechtigungen des Benutzers (beide), Zertifizierung? (nur AD), 
%Zentrale Frage des IAMs: Wer (Login) hat was (Lokale Prozessevents, Netzwerkdaten) wann (Timestamps, generell) wo (Kombination der Elemente) wie (Kombination der Elemente) gemacht?
%Spezifizierte Frage: Welche Kontextinformationen bekomme ich von einer zentralen Benutzerverwaltung, die ich nicht von einem lokalen System bekomme?
%Ansatz: Erstmal Windows anschauen, dann ggf. nach AD suchen
%Temporäres Resultat: AD erstmal weglassen, später einbauen


%Frage: Brauche ich Remote-Access?

\begin{itemize}
\item Router, Rolle: Zugang zum WAN
\item Switch, Rolle: Verbindungen zwischen den Elementen herstellen (Router <-> Switch <-> Firewall...)
\item Firewall, Rolle: Untersuchung des Netzwerkdatenverkehrs 
\item Webserver (Apache), Rolle: Plattform für den Webservice (Shop)
\item Endbenutzer PC, Rolle: Element über das Benutzerzugriffe auf verschiedene Systeme ausgeführt werden
\item Datenbankserver (MSSQL), Rolle: Enthält Businesslogik und kritische Daten
\item IDS (NIDS)?, Rolle (potentiell): Sicherheitssystem mit Event-Daten für das SIEM
\item SIEM, Rolle: Zentrale Log-Verwaltung und Korrelation 
\end{itemize}
\subsection{Applikationen}

%Die Frage ist, welche Applikationen brauche ich?
%#1 Eine Applikation, die auf das Dateisystem eines Elementes zugreift (WebServer)
%#2 Eine Applikation, die über das Netzwerk kommuniziert
%#3 Eine Applikation die Log-Daten erzeugt

%//Nicht zu komplex machen

\subsection{Interaktionen der Systeme (und Anwender)}

% Grobe Übersicht, erster Ansatz
% Anwender <--> Webserver
% WebService <--> Webserver
% Router(GW) <--> Webserver
% Webserver <--> Datenbank Server
% Endbenutzer <--> Datenbank Server
% Endbenutzer <--> Webserver
% Endbenutzer <--> Router(GW)

% Wo überwacht das NIDS (inline?)?
% Was und wie extrahiert das SIEM?

\section{Verwendete Analysemethode}

% Schritt 1: Recherchiere welche Log Dateien existieren
% Schritt 2: Jede Log Datei auf einem System ansehen und Eventformate heraussuchen
% Schritt 3: Events nach Typ kategorisieren
% Schritt 4: 

\section{Resultierende Informationstypen \/ -kategorien}
\subsection{Extraktion und Verarbeitung}
\subsection{Windows}
\subsection{Linux}
\subsection{Beispielapplikation}
\subsection{Netzwerk}
\section{Analyse im Bezug auf Angriffsvektoren}

%------------------KAPITEL : EVALUATION ----------------- %
\chapter{Analyse der Informationsquellen in einem industriellen Produktionsnetzwerk}
\label{cha:Analyse der Informationsquellen in einem industriellen Produktionsnetzwerk}
\section{Beschreibung der Unternehmensarchitektur (Beispiel)}
\subsection{Systeme}
\subsection{Applikationen}
\subsection{Interaktionen der Systeme (und Anwender)}
\section{Verwendete Analysemethode}
\subsection{(Beschreibung der Kategorisierung und/oder wissenschaftliche Grundlage)}
\section{Resultierende Informationstypen / -kategorien}
\subsection{Extraktion und Verarbeitung}
\subsection{Feldbus im Beispiel}
\subsection{Industrial Ethernet im Beispiel}
\section{Analyse im Bezug auf typische Angriffsvektoren (s. Grundlagen)}

%------------------KAPITEL : ERGEBNIS ----------------%
\chapter{Vergleich der Analysen}
\label{cha:Vergleich der Analysen}

In diesem Kapitel wird der Vergleich der Informationsmengen durchgeführt. Als Grundlage für den Vergleich werden die Informationsmengen verwendet, die in den Analysen (Kapitel 3 und 4) ermittelt und in Tabellenform dargestellt sind. Das Ziel dieses Kapitels ist die Erstellung einer Liste, mit Unterschieden zwischen den ermittelten Netzwerkmodellen, die im nächsten Kapitel einer Bewertung unterzogen wird. In diesem Kapitel wird zunächst die verwendete Vergleichsmethode beschrieben, worauf die Anwendung dieser Methode erfolgt und das Ergebnis des Vergleiches dargestellt wird.

\section{Vergleichsmethode}
Die Basisannahme für die Vergleichsmethode ist, dass die Informationsmengen zwischen dem Unternehmens- und dem Industrienetzwerkmodell unterschiedlich sind und Differenzen auftreten. Diese Annahme ergibt sich aus der Tatsache, dass etwa in einem industriellen Netzwerk Systeme verwendet werden, die spezielle Anforderungen des Umfeldes erfüllen müssen und damit bestimmte Limitierungen aufweisen, welche üblicherweise in einem Unternehmensnetzwerk nicht zu finden sind. Dies schließt sowohl Systeme als auch die verwendeten Kommunikationstechnologien ein.
Deswegen werden die Analysen angefertigt. Um den Vergleich unterschiedlicher Systemtypen zu ermöglichen, ist es notwendig, eine gemeinsame Basis zu schaffen. Hier wird dies mit der Unterteilung in System und Kommunikationsprotokolle und auch mit der Unterteilung der Informationen in Kategorien erreicht. \\ 

Diese Basis ermöglicht einen mehrschichtigen Vergleich. In einem ersten Schritt wird untersucht, ob für jede Kategorie jeweils Informationsquellen vorhanden sind. Durch die Informationsquellen des Unternehmensnetzwerkes (hauptsächlich in Form von Protokolldateien) ergeben sich gewisse Unterkategorien, für die die Existenz äquivalenter Informationsquellen im Industrienetzwerkmodell geprüft werden. Im letzten Schritt folgt ein Vergleich dieser äquivalenten Informationsquellen. In diesem werden Unterschiede dieser äquivalenten Informationsquellen herausgearbeitet und Unterschiede notiert.\\


\section{Vergleich}
Die beschriebene Methodik wird zunächst für den Vergleich zwischen Server und SPS angewandt. Das Ziel ist es, die Unterschiede zwischen den beiden Systemtypen zu zeigen. Zusätzlich werden desweiteren die Anwendungen verglichen und Unterschiede aufgezeigt. Es folgt der Vergleich zwischen die Ethernet/IP/TCP-basierte Kommunikation und der Feldbuskommunikation am Beispiel von PROFIBUS. Dadurch werden die grundlegenden Unterschiede herausgestellt. Desweiteren wird PROFINET mit seinen verschiedenen Konfigurationen und Diensten mit den Protokollen der Anwendungsebene verglichen.

\subsection{Systeme}
Im Folgenden werden die Betriebssysteme mit dem SPS-Modell anhand der ermittelten Informationsquellen verglichen.\\

%Ebene 2
%System 
\textbf{Kategorie: System}

Bei der Betrachtung der Kategorie \glqq System\grqq{} sind in allen drei Systemen Informationen über deren Zustand verfügbar. Im Detail protokollieren sowohl die Betriebssysteme als auch die SPS Ereignisse bei Start- und Stop-Vorgängen des Systems. Während jedoch bei den Betriebssystemen die einzelnen Schritte dokumentiert werden, ist die Dokumentation des Betriebszustandwechsels der SPS von STOP zu Anlauf zu RUN limitiert. Die dokumentierten Ereignisse im Diagnosepuffer beschränken sich auf protokollierte Ereignisse der Anlauf-OBs (Standardereignisse) und der damit verbundenen Anzahl der Anlauf-OBs sowie Fehlerereignisse, die während des Anlaufes auftreten. 

Bei Hardwarefehlern werden Fehlerereignisse protokolliert, die durch das Betriebssystem erkannt werden (etwa bei Schnittstellen-Fehlern) sowie Fehlern, die durch die Gerätetreiber an das Betriebssystem gemeldet werden. Daraus ergibt sich, dass die dokumentierten Fehler beschränkt sind durch Treiberdefinition und Betriebssystemkennwerte. Äquivalent werden im Diagnosepuffer der SPS von der CPU erkannte Fehler bzgl. der CPU selbst, des Gerätes, eingebauter Module und angeschlossener Peripheriegeräte dokumentiert. Die Alarme der (Pheripherie-)Geräte sind in der entsprechenden GSD Datei dokumentiert. Desweiteren dokumentiert die SPS auch Kommunikationsfehler zu den Pheripheriegeräten (etwa Kurzschlüsse oder Drahtbrüche).

Änderungen an Betriebssystemkomponenten (inklusive Verzeichnissen, Softwarekomponenten und Dateien) inklusive Sicherheits(sub-)systemen werden in System- bzw. Sicherheitsprotokollen dokumentiert. Erweiterte Ereignisse etwa bzgl. Änderungen an der Windows Registry oder System Daemons (Linux) lassen sich mit zusätzlicher Sicherheitssoftware überwachen. Die S7-1200 dokumentiert äquivalent Änderungen des Betriebssystemzustandes (wie bereits für die Installationen erwähnt). Darüber hinaus können in neueren Versionen weitere Systeminformationen durch die System- und Taktmerker gewonnen werden. Desweiteren werden Sicherheitsereignisse wie etwa die Wiederherstellung einer CPU-Konfiguration, Änderungen an den Schutzeinstellungen der CPU und Änderungen an der Aktualisierungsdatei der Firmware gewonnen. Es ist nach Wissen des Verfassers nicht genau geklärt, ob dabei auch konkrete Änderungen der Konfiguration bzw. der Firmware dokumentiert werden können.

Als letzter Vergleichspunkt werden Fehlermeldungen des Systems betrachtet. Diese Unterkategorie überschneidet sich teilweise mit dem Punkt der Hardwarefehlermeldungen, jedoch werden zudem softwareseitige Fehler bei der Systemausführung mit einbezogen. Diese werden bei den Betriebssystemen in System- und Diagnoseprotokollen festgehalten. Die S7-1200 dokumentiert in diesem Zusammenhang funktionelle Fehler während des Betriebs wie etwa E/A-Fehler durch Unerreichbarkeit oder fehlerhaften Zugriff auf E/A Speicherplätze oder Datenbausteine. Desweiteren werden auch Überschreitungen der maximalen Zykluszeit protokolliert.\\

%Datenablage
\textbf{Kategorie: Datenspeicher}

Änderungen an Dateien werden auf Betriebssystemen durch das verwendete Dateisystem durch Zeitstempel definiert. Neben dem Eigentümer/Ersteller der Datei werden in moderenen Dateisystemen der Zeitpunkt der Erstellung, der letzten Änderung und des letzten Zugriffes in Protokollen festgehalten sowie der Zugriff auf kritische und bestimmte benutzerdefinierte Verzeichnisse. Bei der S7-1200 ist diese Dokumentation deutlich stärker limitiert. So werden im Diagnosepuffer Fehlerereignisse von Zugriffsoperationen (etwa beim Lesen/Schreiben von Datensätzen) dokumentiert. Quellen, die die Protokollierung weiterer Ereignisse belegen, konnten nicht gefunden werden. Das einzige Sicherheitsevent in dieser Kategorie protokolliert, dass eine Änderung von Daten auf einer Speicherkarte der S7-1200 stattgefunden hat.\\

%Prozesse und Dienste
\textbf{Kategorie: Prozesse und Dienste}

Windows und Linux dokumentieren den Lebenszyklus der durchgeführten Prozesse und (Hintergrund-)Dienste. Dabei können detaillierte Informationen wie etwa das Ausführungsverzeichnis, Laufzeit, Speicherverbrauch oder der Elternprozess ermittelt werden. Mithilfe von zusätzlicher Software, wie etwa einer Anti-Malware-Software, lassen sich darüber hinaus verdächtige Prozesse oder Prozessketten erkennen und Sicherheitsereignisse protokollieren. Zusätzlich werden Ereignisse bzgl. Ausführung und Fehler der Hintergrunddienste (Services / \glqq cron jobs\grqq ) protokollieren. Bei der S7-1200 SPS werden für das Aufrufen und Ausführen des Programmes und Unterprogrammen verschieden Organisationsbausteine verwendet, die durch von der CPU erkannte Ereignisse angestoßen und teils protokolliert werden. Für jeden Organisationsbausteintyp wird eine limitierte Liste an Standard- und Fehlerereignissen definiert. Diese Form der Dokumentation schließt die Ausführung und Ausführungspausen von (Unter-)Programmen sowie die zeitgesteuerte Ausführung von Unterprogrammen ein.\\

%Benutzer 
\textbf{Kategorie: Benutzerzugriff}

Sowohl für Windows- als auch Linux-basierte Betriebssysteme werden Zugriffsversuche auf das System in verschiedenen Protokolldateien dokumentiert. Bei Windows werden u.a. auch Loginversuche von lokalen und domänen-abhängigen Logins unterschieden. Im Detail werden fehlgeschlagene als auch erfolgreiche Anmeldungen protokolliert. Desweiteren wird in einem gewissen Rahmen die Nutzung von privilegierten Benutzerrechten sowie Änderungen an den Sicherheitseinstellungen des Systems bzgl. Authentifizierungs- und Authorisierungsregeln festgehalten. Bzgl. der betrachteten SPS sind für den direkten Zugriff auf die CPU bzw. den Zugriff auf den integrierten Webserver verschiedene Sicherheitsereignisse definiert. Die Liste dieser Ereignisse beinhaltet zunächst Ereignisse bzgl. der erfolgreichen oder fehlgeschlagenen Anmeldung an dem Gerät (spezieller: der Online CPU) bzw. dem Webserver (etwa durch Eingabe eines Passwortes). Desweiteren werden weitere Auffälligkeiten dokumentiert wie etwa Zeitablauf einer Online-Verbindung durch Inaktivität, multiple gleichzeitige Zugriffe und Änderungen an den Schutzeinstellungen, speziell des CPU-Schutzes. Die Genauigkeit der Protokollierung dieser Ereignisse (speziell bzgl. der Häufigkeit) kann durch Einstellungen des Diagnosepuffers limitiert werden. Zusätzlich zu dieser Liste könnten weitere Sicherheitsereignisse existieren, zu denen jedoch zum Zeitpunkt des Rechercheabschlusses keine Dokumentation gefunden werden konnte.\\  

%Anwendungen 
\textbf{Kategorie: Anwendungen}

An dieser Stelle soll noch ein kurzer Blick auf die verwendeten Anwendungen sowohl im Unternehmensnetzwerk als auch im Fertigungsnetzwerk geworfen werden. Alle drei Anwendungen (Apache Webserver, Microsoft SQL Server, Siemens WinCC) setzen auf Windows oder Linux auf. Der Apache Webserver und Microsoft SQL Server dokumentieren beide verschiedene Aspekte in verschiedenen Protokolldateien. Dabei werden Ereignisse teilweise in Systemprotokolle integriert. Zu den dokumentierten Bereichen zählen u.a. die Dokumentation von Benutzerzugriffen, Datenanforderungen über die jeweilige Schnittstelle, Senden von Befehlen und Anwendungsfehler. Damit werden verschiedene Bereiche überwacht und protokolliert, die genutzt werden können, um den Zugriff auf gespeicherte Daten in der jeweiligen Anwendung nachzuvollziehen.
Für WinCC ergeben sich ähnliche Fähigkeiten. Der Begriff \glqq Protokollierung\grqq{} wird in WinCC als Ausdruck für die Möglichkeit verwendet, verschiedene Datensätze auszudrucken. Neben der Dokumentation von ein- und ausgehenden Prozessdaten, Projektdaten und Änderungen an diesen Projektdaten kann mit Hilfe der WinCC/Audit Option eine Dokumentation der Benutzeraktionen erreicht werden. Dabei werden sämtliche Aktionen des Benutzers vom erfolgreichen/fehlerhaften Login (inklusive Methode) über die Ausführung von Aktionen innerhalb des Projektes bis zur Logout-Methode nachvollzogen werden.


\subsection{Kommunikationsprotokolle}
Für den Vergleich der Kommunikationsprotokolle werden die Kategorien \glqq Header\grqq{} und \glqq Verhalten\grqq{} verwendet. Da jedoch die Protokolle unterschiedliche Funktionalitäten erfüllen, wird im Folgenden darüber hinaus zwischen Kommunikations- und Anwendungsprotokollen unterschieden. Die Kommunikationsprotokolle ermöglichen die Datenübertragung und die Anwendungsprotokolle setzen darauf auf.

Für die Unterteilung werden daher zunächst PROFIBUS und die Kombination von Ethernet, IP und TCP/UDP verglichen. Darauf folgend werden die Anwendungsprotokolle des Unternehmensnetzwerkes und die Gemeinsamkeiten mit Informationen über Anwendungen bzw. Diensten von PROFIBUS und PROFINET jeweils verglichen. \\


\subsubsection{Vergleich von PROFIBUS und Ethernet-basierter Kommunikation}
\textbf{Header}

Bei dem Vergleich von PROFIBUS und Ethernet-basierter Kommunikation sind zunächst die Gemeinsamkeiten zu nennen. Diese bestehen in einer Sender- und Empfängeradresse der Nachricht, einer Kennung der Nachricht für einen bestimmten Typ sowie einer Kennung für einen bestimmten Dienst bzw. Service-Typ. 

Der fundamentale Unterschied ergibt sich aus der Menge der möglichen Kodierungen bzw. Detailtiefe der Informationen über die verschickte Nachricht selbst. So finden sich Daten, wie etwa die Gültigskeitsdauer einer Nachricht oder Informationen über die Länge, in jeder Ethernet-basierten Nachricht, während die Gültigkeitsdauer überhaupt nicht und die Länge der Nachricht nur für ein bestimmtes Nachrichtenformat in der PROFIBUS-Kommunikation versendet wird. Bei PROFIBUS werden teilweise Paritäts-Bits oder CRC-Daten innerhalb der Nutzdaten mitgesendet. 

Eine weitere Gemeinsamkeit ergibt sich aus der Übermittlung von Port-Daten bzw. SAP (Service Access Point) Daten. Diese Informationen definieren äquivalent den Zugangspunkt einer Anwendung bzw. eines Dienstes am System.\\

\textbf{Informationen durch Verhaltensanalyse}

Bei der Betrachtung der Verhaltens der Kommunikationssysteme im laufenden Betrieb lassen sich Gemeinsamkeiten finden in Form der Nachvollziehbarkeit bei der Eingliederung von neuen Teilnehmern. In der Ethernet-basierten Kommunikation wird von verschiedenen Protokollen wie etwa dem Adress Resolution Protocol (ARP) eine Zuweisungsliste am Switch bzw. Router des jeweiligen Netzwerkes erzeugt. Diese Aktivität kann dokumentiert werden. Bei der PROFIBUS Kommunikation kann durch das FDL das Hinzufügen bzw. Entfernen von Kommunikationsteilnehmern durch Mitschneiden der Kommunikation ermittelt werden und durch eine Option eines PROFIBUS-Master-Gerätes lassen sich die Kommunikationsteilnehmer im Feldbus durch Abfrage ermitteln. 

Bei der Kommunikation zwischen Geräten lassen sich zusätzliche Daten über den Kommunikationsaufbau bei der Verwendung von TCP gewinnen. So wird bei der Kommunikation zwischen Sender und Empfänger eine Sequenz- bzw. Bestätigungsnummer verschickt. Dieser Punkt ist erwähnenswert, da bei PROFIBUS ebenfalls Bestätigungen für den Empfang von Nachrichten durch die PROFIBUS-Slaves gesendet werden, jedoch die Bestätigung allgemein gehalten wird und keine Daten über eine bestimmte Nachricht in der Bestätigung vorhanden ist. Bei der Verwendung von UDP werden nach der Protokollspezifikation keine Bestätigungen versendet und keine Sequenznummern übertragen.

Bei der Verwendung von TCP durch den TCP-Handshake und das Austauschen von Informationen für die Herstellung einer Verbindung. Die Aufzeichnung dieses Verhaltens ergibt Informationen über den konkreten Verbindungsaufbau zwischen zwei Kommunikationsteilnehmern. Bei PROFIBUS lassen sich diese Informationen aus der Analyse des Nachrichtenverlaufs und die Ermittlung der jeweiligen Empfängeradresse ermitteln. Da genau ein Kommunikationsteilnehmer zur Zeit auf dem Kommunikationsbus senden darf und ein PROFIBUS-Slave-Gerät eine temporäre Sendeerlaubnis vom PROFIBUS-Master für das Senden von Prozessdaten bzw. Bestätigungen erhält, lässt sich auf diese Art rudimentär ein ähnliches Verhalten nachvollziehen. Abhängig von dem verwendeten Übertragungsdienst kann jedoch nicht ermittelt werden, ob die Nachricht erhalten wurde, wenn eine Quittierung nicht verlangt wird. Hier lässt sich eine Parallele zu UDP ziehen.

Fehler innerhalb der Kommunikation lassen sich durch Fehlermeldungen der Kommunikationsteilnehmer in Ethernet-basierten Netzwerken nachvollziehen (etwa in Form von ICMP-Nachrichten). Bei PROFIBUS-Systemen können durch Dienste des FDMA-Protokolls Funktionsfehlernachrichten verschickt werden. \\


\subsubsection{PROFIBUS-Dienste und Anwendungsprotokolle des TCP/IP Stacks}
\textbf{Gemeinsamkeiten der Kommunikationsprotokolle}

Bei erster Betrachtung verbindet die ausgesuchten Protokolle ihre grundlegende Eigenschaft der Client-Server-Anwendung, d.h. jedes Protokoll dient dazu, den Lebenszyklus der Verbindung von einem Client-Gerät zu einem Server-Gerät zu verwalten und den Austausch zu kontrollieren. 

An dieser Stelle hebt sich HTTP in der Form von den restlichen Netzwerkprotokollen derart ab, dass es für die Verbindungsverwaltung und -abwicklung nicht mehrere Subprotokolle verwendet. SSH, RDP und MSSQL benutzen Protokollsuiten mit dedizierten Protokollen für verschiedene Schritte. So werden Protokolle verwendet, um die Verbindung zu initialisieren oder die Verbindungskonfiguration zu regeln. Ein weiteres Subprotokoll dient dann der Ausführung der eigentlichen Anwendung durch den Nutzer. Auffallend sind auch die vorhanden Authentisierungsmechanismen, die durch ein eigenes Subprotokoll durchgeführt werden.

Desweiteren lassen sich die Protokolle zunächst in zwei Arten unterteilen. Die Protokolle der MSSQL Suite und das HTTP dienen der Anfrage von Ressourcen und Diensten an den Server und den Erhalt einer Bestätigung und/oder Ressource. SSH und RDP hingegen dienen dem Zugriff auf das Serversystem. 

%Header
Bei der Betrachtung der Header der versendeten Nachrichten kann von jedem Protokoll Name und -version sowie Nachrichtentyp ermittelt werden. Darüber hinaus werden weitere Informationen über die gewählte Verbindung mitgesendet, wie etwa die verwendeten Verschlüsselungsmethoden (für  HTTP über TLS (Schicht 5 im OSI Modell))  und Informationen über die Nachricht in Form von Längendaten und weiteren Daten über den versendeten Payload. Alle Protokolle haben die Möglichkeit, verschiedene Verschlüsselungs- und Authentifizierungsmechanismen einzusetzen.

%Verhalten
Der Aufbau der Verbindung für SSH, RDP und MSSQL erfolgt in festgelegten Sequenzen. So wird eine initiale Verbindung erzeugt, die Authentizität des Servers geprüft, Konfigurationsdaten für die Verbindung ausgetauscht und eine Authentizitätsprüfung des Benutzers durchgeführt. Für HTTP gilt dies nicht, da die Kommunikation zustandslos verläuft, d.h. das Anfragen und Antworten ausgetauscht werden ohne Kenntnis des vorhergegangenen Austausches.\\

\textbf{Vergleich mit PROFIBUS-Diensten}

Für den Vergleich der Anwendungsprotokolle werden im Folgenden PROFIBUS DP-V0 und DP-V1 betrachtet. Diese setzen auf der Kommunikationsfunktionalität auf, die durch das FDL-Protokoll bereitgestellt wird. Dies bildet einen Äquivalenzpunkt zu der Art und Weise, wie die Anwendungsprotokolle auf der Ethernet-basierten Kommunikation aufsetzen. 

%Verhalten des Systems
Als ersten Vergleichspunkt wird der Verbindungsaufbau betrachtet. Ähnlich wie bei einem Client-Server Modell werden über das Protokoll Anfragen geschickt und und Antworten erhalten. Für die Initialisierung eines PROFIBUS-DP-Slave-Gerätes wird, wie bei SSH, RDP und MSSQL, eine Verbindung hergestellt und Konfigurationsdaten gesendet. Im Unterschied ist der Austausch der Daten allerdings einseitig, da die Informationen nur von dem Master-Gerät an das Slave-Gerät gesendet und der Empfang durch den Slave quittiert wird. Desweiteren findet keine Prüfung der Authentizität des Masters statt. Zwar wird eine Identifikationsnummer gesendet, anhand derer das Slave-Gerät prüft, ob die über den gemeinsamen Datenbus gesendete Konfiguration auf dem Gerät angewendet werden soll. Darüber hinaus wird allerdings keine weitere Prüfung durchgeführt. Es werden zudem keine Verschlüsselungsmethoden ausgetauscht. Die Definition der Parameter wird über die Daten der GSD-Datei definiert, die dem Master-Gerät zur Verfügung stehen. In Folge dessen wird die Konfiguration zwar innerhalb der Sequenz geprüft, es findet jedoch im Vergleich zu SSH, RDP und MSSQL keine Aushandlung der gemeinsamen Konfiguration für die Verbindung statt.

%Header
Eine Gemeinsamkeit der Headerdaten zwischen PROFIBUS DP-V0/V1 und den Anwendungsprotokollen besteht in der Information über die Länge der Anfrage/Antwort-Nachricht sowie des Nachrichtentyps. Desweiteren werden bei der DP-V1 Kommunikation äquivalent zu HTTP Informationen über die angefragte Ressource übertragen. 
Bei Verwendung der synchronen Kommunikation mit DP-V0 werden bei der Initialisierung im Zuge der Parametrisierung Daten über die verwendbaren Alarmtypen, die Gruppenzugehörigkeit für die Anwendung der Broadcast-Funktionalität durch einen Master der Klasse 1 sowie weitere Benutzerparameter übertragen. 
Im Vergleich zu den SSH, RDP und MSSQL werden diese Konfigurationsdaten unverschlüsselt übertragen. 

Eine weitere Gemeinsamkeit besteht in der Existenz eines spezifizierten Fehler-Nachrichtentyps, der einen Fehler der Kommunikation beschreibt. Im Unterschied zu den Anwendungsprotokollen des Unternehmensnetzwerk variiert jedoch die Anzahl der Typen abhängig von dem Kommunikationspartner. Während dies eingeschränkt auch für unterschiedliche Servertypen zutrifft, wird jedoch eine Grundmenge über das Protokoll definiert. Darin liegt der entscheidende Unterschied.\\


\textbf{Vergleich PROFINET und PROFIBUS}

Als Grundlage basiert PROFINET zunächst auf dem Ethernet-Standard. Abhängig von dem gewählten Kommunikationskanal werden jedoch UDP/IP bzw. TCP/IP verwendet oder nur auf Basis von Ethernet. In diesem Fall besteht eine Nachricht aus einem Rahmen, der eine PROFIBUS-Nachricht zwischen den Ethernet-Headern und einem CRC-Feld einkapselt. Auf der Anwendungsebene betrachtet man die PROFINET Dienste und Protokolle \citep{PROFINET3}.
Die Dienste beziehen sich prinzipiell auf die bereits verglichenen Read/Write-Dienste von PROFIBUS. Zusätzlich wird durch die verwendete Kommunikationbeziehung (CM, IO data oder Alarm) eine zusätzliche Information über den allgemeinen Zweck der Nachricht gewonnen. Diese Beziehung und die zusätzlichen Protokolle DCP und LLDP (vgl. Industrieanalyse/PROFINET) erlauben eine genauere Eingrenzung des Nachrichtenzweckes, äquivalent zu den Subprotokollen von SSH, RDP und MSSQL. Die verwendeten Protokolle ermöglichen es im Vergleich zu einer PROFIBUS-Kommunikation, zusätzliche Informationen über die Identität eines Kommunikationsteilnehmers zu erhalten sowie Prozesse der Adressvergabe nachzuvollziehen. Dies ermöglicht die genauere Segmentierung des Nachrichtenverkehrs und in der Folge zusätzliche Informationen über den Nachrichtenaustausch im Netzwerk.

\section{Ergebnis}
Als Ergebnis des Vergleiches ist festzustellen, dass grundlegende Ähnlichkeiten der Informationsmengen zu erkennen sind. Dies ist ein natürlicher Schluss aus dem Zweck der Anwendung. Es ergeben sich jedoch auch Unterschiede. \\

Zunächst kann sowohl für die verglichenen Systeme als auch Kommunikationsprotokolle festgehalten werden, dass die Elemente des industriellen Netzwerkes nach dem vorliegenden Vergleich eine eingeschränktere Informationstiefe vorweisen. Im Speziellen ergeben sich bei der S7-1200 Limitierungen in der Dokumentation von Ereignissen im Diagnosepuffer, die hauptsächlich funktionelle Fehlermeldungen und Systemereignisse zeigen. Dies ist besonders präsent in der Dokumentation von Parameteränderungen. Die potentielle Einschränkung der Protokollierung von Benutzerzugriffen (im Sinne der Häufigkeit und Vielfalt der Ereignisse), basierend auf der Auslastung des Diagnosepuffers, schränkt die Möglichkeit der Dokumentation von Änderungen ebenfalls ein. \\

Bei dem Vergleich der Kommunikationsprotokolle ergaben sich die folgenden Unterschiede: Zum einen konnte keine Unterstützung von Integritätsprüfung der Kommunikationspartner bei der Verwendung des PROFIBUS-Protokoll festgestellt werden. Darüber hinaus ergeben sich Unterschiede in der Art der Festlegung verfügbarer Nachrichtentypen für die Kommunikation zwischen zwei Geräten auf der Anwendungsebene.



\chapter{Bewertung der Informationslücken}
\label{cha:Bewertung der Informationslücken}
\section{Bewertungsschema}
\section{Bewertung}
\section{Beschreibung existierender wissenschaftlichen Lösungsansätze}

\chapter{Lösungsansatz zur Schlie\ss ung der fokussierten Informationslücke}
In diesem Kapitel wird die Informationlücke der Protokollierung von Änderungen an Datenbausteinen und Parametern untersucht und ein Ansatz für die Schließung dieser Informationslücke beschrieben.

\section{Ansatzentwicklung und Szenarien}
Das Ziel des Ansatzes ist es, eine Möglichkeit zu finden, Änderungen an E/A-Parametern sowie Bausteinen zu dokumentieren. 

\subsection{Annahmen und Szenarien}
Zu diesem Zweck werden zunächst die folgenden Annahmen definiert:
\begin{itemize}
\item Die Protokollierung der Änderungen an Parametern und Bausteinen ist aktuell nicht möglich
\item Für eine skalierbare Lösung ist eine dauerhafte Überwachung nach aktuellem Stand nicht möglich
\item Für die Komprimittierung einer SPS ist die Infizierung eines Servers mit Kommunikationsschnittstelle zu der Malware notwendig
\end{itemize}

Diese Annahmen werden im Folgenden weiter ausgeführt und begründet.

\subsubsection{Keine Änderungsprotokollierung auf der SPS}
Für diese Hypothese sprechen mehrere Faktoren. Zum einen sind die Werte von Prozessdaten und der zugehörigen Speicherplätze sehr volatil. 
Es wird angenommen, dass sich aus der Zykluszeit der Ausführung eines Rezeptes, der zugehörigen Änderung der Sensorwerte und der daraus resultierenden Anpassung der Ausgangswerte der verschiedenen angeschlossenen Feldgeräte eine hohe Wertefluktuation ergibt. 
Abhängig von der Größe des Ladespeichers und der Verfügbarkeit von zusätzlichem Speicherplatz durch eine Speicherkarte ist zudem die Kapazität des Ladespeichers und damit die Anzahl der möglichen Einträge begrenzt. 
Ein weiteres Argument ist, dass die Protokollierung von Prozessdaten auf einer SPS durch den Programmierer des Programmes eingepflegt werden müssen. Abhängig von der Zykluszeit und der notwendigen Operationen kann die Anzahl an möglichen Schreiboperationen stark begrenzt sein. Laut dem Systemhandbuch der im Modell verwendeten SPS S7-1200 kann bspw. nur die Erstellung einer Datenprotokolldatei einige Zeit in Anspruch nehmen. Desweiteren ist die Limitierung einer Datenprotokolldatei auf 255 Einträge eine weitere Begrenzung.

\subsubsection{Keine dauerhafte Überwachung bei skalierbarer Lösung}
Diese Hypothese wird aufgestellt basierend auf der Anzahl der zu überwachenden Geräte an einem Feldbus, z.B. PROFIBUS. Durch den gemeinsamen Datenbus steht für ein potentielles Gerät, welches Datenanfragen an SPS-Teilnehmer schickt, nur ein begrenzter Zeitrahmen zur Verfügung. Die Erhöhung des Zeitrahmens könnte die Prozessdurchführung einschränken bzw. nicht die dafür notwendigen Reaktionszeiten erreichen, die für die Prozessdurchführung notwendig sind. 

\subsubsection{Infizierung eines Servers für SPS-Manipulation notwendig}
Diese Hypothese leitet sich aus der Vorgehensweise des Stuxnet-Angriffes ab. Die Manipulation der Kommunikationsschnittstelle kann bspw. durch reinen Schadcode auf der SPS nur schwierig verhindert oder verborgen werden. Zu diesem Zweck müssten Anwenderprogramm und Firmware derart manipuliert werden, dass die Korrektur der eigentlichen Werte, wie sie bei Stuxnet durch die Rootkit-Funktionalität durchgeführt werden, direkt auf der SPS mit der zur Verfügung stehenden, unter Umständen sehr geringen, Rechenleistung durchgeführt werden. Die Infizierung der SPS ist zudem ohne eine Infizierung eines Servers mit einer Verbindung zu der betroffenen SPS nur durch einen Insider möglich, dessen Aktivitäten ohne spezielle Kenntnisse, etwa durch das WinCC-SCADA-System, protokolliert werden.

\subsubsection{Szenarien}
Aus diesen Hypothesen ergeben sich die folgenden potentielle Angriffsszenarien für die Manipulation des SPS-Verhaltens:

\begin{itemize}
\item Manuelle Manipulation von Parametern durch einen Insider 
\begin{itemize}
\item über ein Hardware-Programmiergerät (erfordert physikalischen Zugriff)
\item über einen angeschlossenen Server mit installierter Programmiersoftware (etwa WinCC)
\end{itemize}
\item Manipulation per Malware
\begin{itemize}
\item Autonome Malware wie am Beispiel Stuxnet
\item Manipulierbare Malware, die etwa eine eigene Kommunikationsverbindung mit Hilfe der Serverschnittstellen aufbauen kann (erfordert Zugriff des Angreifers entweder direkt oder indirekt über eine Hintertür)
\end{itemize}
\end{itemize}

Auf der Grundlage der beschriebenen Hypothesen und Szenarien lassen sich zudem weitere Schlussfolgerungen für Anforderungen an einen möglichen Lösungsansatz generieren. 

Basierend auf der weiteren Hypothese, dass eine dauerhafte Überwachung nicht stattfinden kann, lässt sich die Schlussfolgerung ziehen, dass die Erhebung und Prüfung der SPS-Daten zu bestimmten Zeitpunkten, etwa durch ein Ereignis, ausgelöst werden muss. Dieses Ereignis kann unterschiedlichen Szenarien entstammen. 
In Bezug auf das Stuxnet-Beispiel könnten sich z.B. Hinweise bei einem Abgleich der Prozessdaten der gesendeten/empfangenen Nachrichten ergeben. In diesem Falle wird ein Unterschied bei gesendeten Frequenzdaten festgestellt, die beim Aufruf der Write-Methode durch Stuxnet manipuliert werden. 
Ein weiterer Hinweis kann aus einem Abgleich der Anzahl der durch einen Benutzer angeordneten und der an die SPS gesendeten Anfragen resultieren. Hier wird im Zuge des Abgleiches der Anzahl der Nachrichten eine Differenz entstehen durch die regelmäßigen Anfragen der Stuxnet-Malware bzgl. des manipulierten Datenbausteines. 

Eine weitere Schlussfolgerung ergibt sich aus der limitierten Protokollierungskapazität. Durch diese Limitierungen ist es notwendig, die Protokollierung der Nachrichten etwa durch ein dediziertes Element wie einen weiteren Server als DP-V1-Master der Klasse 2 durchführen zu lassen. 

\section{Beschreibung des Lösungsansatzes}
Der folgende Lösungsansatz wird basierend auf den vorher geschilderten Annahmen und Schlussfolgerungen erstellt. Eine Übersicht über Datenfluss und Elemente des Ansatzes sind in der folgenden Abbildung zu sehen.
\todoForm{Lösungsansatz: Bild für Ansatzkonzept einfügen}

Für den Lösungsansatz wird angenommen, dass eine Möglichkeit besteht, die Kommunikation zwischen SCADA-Server und SPS über ein Switch zu leiten oder über eine andere Möglichkeit den Nachrichtenverkehr durch ein externes Gerät zu überwachen. Aus den beschriebenen Hypothesen und Schlussfolgerungen ergibt sich die Notwendigkeit, dass eine unabhängige Datenquelle vorhanden sein muss, auf deren Basis unabhängig Informationen über den Datenverkehr erheben lassen. Diese Notwendigkeit besteht, da die Integrität der Kommunikationsdaten auf dem SCADA-Server bei einer Infizierung nicht gewährleistet werden kann. 

Mit Hilfe dieses \glqq Netzwerkmonitors\grqq{} kann der Datenaustausch protokolliert werden und an einen weiteren Server weitergegeben werden, der etwa als Zwischenstation für ein SIEM-System oder eine andere Sicherheitsanwendung verwendet werden kann. Dieser Server wird im Folgenden als \glqq Agenten-Server\grqq{} bezeichnet. 
Die Aufgabe dieses Servers ist es, Datenprotokolle der SCADA-Anwendung und Datenprotokolle des Netzwerkmonitors zu vergleichen und Unterschiede in den Datensätzen festzustellen. Ist eine solche Differenz festgestellt, wird basierend auf den Details des Unterschiedes (etwa Unterschied zwischen den Prozessdaten oder Unterschiede in der Nachrichtenanzahl) die Aufgabe der Prüfung der entsprechenden SPS an ein unabhängiges Element weitergeleitet. 
Dies wird durch einen weiteren Server oder Industrie-PC mit einer Kommunikationsschnittstelle zu der entsprechenden SPS umgesetzt. Dieses Element wird im folgenden als \glqq Prüfer\grqq{} bezeichnet. Die Aufgabe dieses Elementes ist es, einen bestimmten Datensatz aus der SPS auszulesen und diesen an den Agenten-Server zu senden. 
Der Agenten-Server vergleicht den Datensatz mit den zugehörigen Werten des SCADA-Systems. Diese zusätzliche Prüfung ist notwendig, um sicherzugehen, dass der übertragene Wert des Netzwerkmonitors dem auf der SPS gespeicherten Wert entspricht und ein Unterschied tatsächlich vorliegt. Wird ein Unterschied festgestellt, ergibt sich daraus die Information, dass eine Änderung der SPS-Daten ohne Dokumentation im SCADA-System vorgenommen wurde.

\subsection{Limitierungen des Ansatzes}
Der beschriebene Ansatz ermöglicht eine grundlegende Prüfung der SPS-Daten auf Änderungen, jedoch ist die Problematik der Datenmanipulation auf dem System vor der Übertragung der Daten an das Prüfer-System nicht auszuschließen. Eine weitere Problematik könnte im Bezug auf die Funktionsweise der SPS und die Dauer der Übertragung auftreten. So kann etwa bei dem Vergleich der Datensätze basierend auf dem Zeitstempel ein Unterschied der Daten bestehen. Aus diesem Grund funktioniert der beschriebene Ansatz nur annäherungsweise.

\section{Demonstration des Lösungsansatzes}
Für die Demonstration des beschriebenen Ansatzes wurde eine abstrakte Simulation erstellt. Da der Erwerb der notwendigen Software-Lizenzen außerhalb des finanziell sinnvollen Rahmens liegt, wurden für die Demonstration notwendige Charakteristiken der Elemente in Form von Java-basierten Klassen umgesetzt.

Dafür wurden die folgenden Software-Pakete erstellt:
\begin{itemize}
\item Kommunikationspakete: Dieses Paket definiert notwendige Kommunikationsroutinen und Nachrichtenformate (Orientierung an PROFINET-Formaten).
\item SPS: Dieses Pakete enthält repräsentativ anfragbare Klassen (Bausteine und E/A-Speicher), eine \glqq Firmware\grqq{} für die Ausführung von der Bausteinklassen sowie eine Verwaltung zu Kommunikationszwecken.
\item SCADA\_Server: Dieses Paket nutzt die definierten Kommunikationsroutinen für die Kommunikation mit der SPS und dem Agenten-Server. Zusätzlich werden die Daten der SPS während der Anwendung gespeichert.
\item Netzwerkmonitor: Diese Komponente dokumentiert die Kommunikation zwischen SCADA-Server und SPS. 
\item Prüfer: Diese Komponente enthält eine Logik, die abhängig von der empfangen Nachricht des Agentenservers eine bestimmte Komponente der SPS anfragt.
\item Agenten-Server: Verwaltung des Anfrageprozesses, ausgelöst durch den Netzwerkmonitor. Enthält Vergleichsroutinen für den Vergleich der empfangenen Daten.
\end{itemize} 

Das folgende Schaubild zeigt die Struktur der Demonstration:
\todoForm{Lösungsansatz: Schaubild des Demonstrationsnetzwerkes einfügen}

Die versendeten Kommunikationspakete haben die folgende Struktur:
\todoForm{Lösungsansatz: Bild des Nachrichtenformates einfügen}

Der Prozessablauf des SPS-Paketes simuliert in einem einfachen Beispiel den Prozess einer SPS. Dabei wird in einer definierten Zykluszeit ein definiertes \glqq Anwenderprogramm\grqq{} ausgeführt, welches die Werte im E/A-Array durch simple mathematische Operationen modifiziert. Empfängt die SPS eine Nachricht des WinCC-Servers werden die Anpassungen der Werte vorgenommen. Das Modell orientiert sich dabei an der echten Adressierung einer SPS. So werden E/A-Werte in verschiedenen Arrays gespeichert, die durch einen Zahlenwert (Slot-Number) adressiert werden können.

Der SCADA-Server sendet in regelmäßigen Abständen Änderungen für E/A-Daten an die SPS. Die Modifikation der Daten wird dabei durch eine mathematische Funktion zwischen den Write-Nachrichten verändert. Ein Zufallszahlengenerator berechnet dabei die Chance einer Negation des entsprechenden Wertes. Diese Modifikation stellt den Manipulationsversuch einer Malware da. 

Da die Simulation einer weiteren Komponente in Form eines Switches weiteren Aufwand erfordert hätte, wird bei der Anwendung der Kommunikationsroutinen die Nachricht jeweils an die SPS bzw. den SCADA-Server und jeweils auch an den Nachrichtenmonitor gesendet. Dieser prüft den enthalten Wert auf das Vorzeichen. Bei einem negativen Wert wird eine Nachricht an den Agentenserver versendet, die die Werte der verdächtigen Nachricht enthält.

Der Agentenserver lauscht auf dem Socket. Sobald eine Nachricht des Netzwerkmonitors eingeht, wird eine Nachricht mit der entsprechenden Speicherstelle an den Prüfer gesendet. Dieser fragt bei der SPS den Wert der entsprechenden Quelle ab und leitet diesen Wert an den Agenten-Server zurück. 

Der Agentenserver fordert den aktuellen Wert der Speicherstelle von dem SCADA-System an und vergleicht die Werte. 

\todoForm{LösungsansatzDemo: Falls noch Zeit ist, Wireshark Bild einfügen}
Der vollständige Quellcode ist auf der DVD im Einband der Masterarbeit hinterlegt.

\section{Zusammenfassung}
In diesem Kapitel wurden Hypothesen und Schlussfolgerungen bzgl. eines möglichen Ansatzes für die Protokollierung von Datenänderung auf einer SPS ausgeführt. Basierend auf diesen Kriterien wurden ein Lösungsansatz und eine demonstrativen Umsetzung des Ansatzes beschrieben.
%----------------- KAPITEL : FAZIT  ----------------- %	

\chapter{Fazit}
\label{cha:Fazit}



%----------------- KAPITEL : BEISPIELE  ----------------- %	
	
\pagenumbering{roman}	

% ----- Referenzen ----- %
%\begin{thebibliography}{6}

%\end{thebibliography}

%\addcontentsline{toc}{chapter}{Literaturverzeichnis}
\appendix

 \bibliographystyle{plainnat}          % Style einstellen
    \bibliography{MeineQuellen}             % eine .bib Datei (Textdatei)
 \addcontentsline{toc}{chapter}{Literaturverzeichnis}
 
%Literaturverzeichnis
% ----- Abbildungen ----- %
%\addcontentsline{toc}{section}{Abbildungsverzeichnis} % falls in Inhalsverzeichnis
\listoffigures
\addcontentsline{toc}{chapter}{Abbildungsverzeichnis}
	

% ----- Tabellen----- %
% \addcontentsline{toc}{section}{Tabellenverzeichnis}  % falls in Inhalsverzeichnis
%\fancyhead[L]{Abbildungsverzeichnis / Abkürzungsverzeichnis} 

%\addcontentsline{toc}{chapter}{Abkürzungsverzeichnis}
\phantomsection \addcontentsline{toc}{chapter}{Abkürzungsverzeichnis}
\renewcommand\refname{Abkürzungsverzeichnis} 
\chapter*{Abkürzungsverzeichnis}

\begin{acronym}[Bash]
 \acro{VM}{Virtuelle Maschine}
 \acro{KDE}{K Desktop Environment}
 \acro{SQL}{Structured Query Language}
 \acro{Bash}{Bourne-again shell}
 \acro{JDK}{Java Development Kit}
 \acro{SPS}{Speicherprogrammierbare Steuerung}
\end{acronym}

%Kopfzeile links
%\listoftables
	
\end{document}